%
% main.tex -- Paper zum Thema Laplace
%
% (c) 2020 Autor, OST Ostschweizer Fachhochschule
%
% !TEX root = ../../buch.tex
% !TEX encoding = UTF-8
%
\chapter{Die Laplace-Inversion: Vom Bromwich-Integral zur numerischen Talbot-Kontur \label{chapter:laplace}}
\kopflinks{Die Laplace-Inversion: Vom Bromwich-Integral zur numerischen Talbot-Kontur}
\begin{refsection}
\chapterauthor{Simon Widmer und Noel Karlsson}

Diese Arbeit widmet sich der Inversion der Laplace-Transformation jenseits gängiger Korrespondenztabellen. 
Im Zentrum steht dabei das komplexe Bromwich-Integral. 
Der Text beleuchtet dessen analytische Auswertung mittels Pfaddeformation und Residuensatz und zeigt, wie sich anspruchsvolle Singularitäten -- etwa Polketten oder Verzweigungsschnitte aus der Wärmeleitungsgleichung -- elegant handhaben lassen. 
Den Abschluss bildet die numerische Inversion: Die Deformation auf die Talbot-Kontur überführt das oszillierende Integral in eine rasch abklingende Form und ermöglicht eine hocheffiziente numerische Auswertung. 
So schlägt das Paper die Brücke von der reinen Funktionentheorie zur algorithmischen Praxis.


\section{Die Laplace-Transformation}
\label{laplace:section:intro}

Die Laplace-Transformation ist ein unverzichtbares Werkzeug in der angewandten Mathematik und den Ingenieurwissenschaften.
Ihre wahre Stärke entfaltet sie bei der Lösung linearer partieller Differentialgleichungen, wie etwa der Wärmeleitungsgleichung.
Durch die Transformation werden zeitliche Ableitungen in einfache algebraische Multiplikationen überführt.
Was im Zeitbereich eine komplexe Differentialgleichung ist, wird im Bildbereich zu einer deutlich einfacheren algebraischen Gleichung.

In der ingenieurmässigen und physikalischen Praxis berechnet man die Laplace-Transformation $\mathcal{L}\{f(t)\}$ fast nie von Hand über das definierende Integral.
Stattdessen nutzt man \emph{Korrespondenztabellen}, die vorberechnete Paare von Zeitfunktionen $f(t)$ und zugehörigen Bildfunktionen $F(s)$ auflisten.
Hat man die algebraische Gleichung im Bildbereich gelöst, reicht oft ein einfacher Blick in die Tabelle, um die physikalische Lösung $f(t)$ im Zeitbereich zu finden.
Einige der wichtigsten Korrespondenzpaare sind in Tabelle~\ref{laplace:tab:korrespondenzen} zusammengefasst.

\begin{table}[htbp]
    \centering
    % Zeilenabstand etwas vergrössern, damit die Brüche gut lesbar sind
    \renewcommand{\arraystretch}{1.8} 
    
    % Die Option 'hbox' zwingt die Box, exakt so breit wie ihr Inhalt (die Tabelle) zu sein
    \begin{tcolorbox}[colback=white, colframe=black, arc=3pt, boxrule=0.8pt, hbox]
        % Spalten: l (links), r (rechts), c (zentriert), l (links)
        \begin{tabular}{l r c l}
            \textbf{Beschreibung} & \textbf{Zeitbereich} $f(t)$ & $\laplace$ & \textbf{Bildbereich} $F(s)$ \\
            \midrule
            Potenzfunktion & $H(t) \cdot t^n$ & $\laplace$ & $\frac{n!}{s^{n+1}}$ \\[1ex]
            Exponentialfunktion & $e^{at}$ & $\laplace$ & $\frac{1}{s-a}$ \\[1ex]
            Dirac-Stoss (Impuls) & $\delta(t)$ & $\laplace$ & $1$ \\[1ex]
            Ableitung & $f'(t)$ & $\laplace$ & $sF(s) - f(0)$ \\
        \end{tabular}
    \end{tcolorbox}
    
    \vspace{0.2cm}
    \caption{Wichtige Laplace-Korrespondenzen; $H(t)$ bezeichnet die Heaviside-Sprungfunktion.}
    \label{laplace:tab:korrespondenzen}
\end{table}

Das Konzept der Tabellen stösst jedoch schnell an seine Grenzen.
Was passiert bei komplexen physikalischen Systemen, wie etwa der instationären Wärmeleitung in einem Stab?
Die sich dort ergebenden Gleichungen liefern Bildfunktionen $F(s)$, die so kompliziert sind -- oft behaftet mit Wurzeln oder hyperbolischen Funktionen --, dass sie in keiner Tabelle zu finden sind.
In genau diesen Fällen, wenn das einfache Nachschlagen versagt, müssen wir die Transformation mathematisch streng umkehren.
Wir sind gezwungen, die komplexe Zahlenebene zu betreten und das sogenannte \emph{Bromwich-Integral} direkt auszuwerten.

Dieses Paper widmet sich exakt diesem Problem:
Wir leiten die fundamentale Inversionsformel her, untersuchen, wie man komplexe Singularitäten wie Verzweigungsschnitte und unendlich viele Pole mathematisch sicher umgeht, und schliessen mit der hochmodernen numerischen Inversion über die Talbot-Kontur ab.

\section{Herleitung aus der Fourier-Transformation}
\label{laplace:subsection:herleitung}

Um die Natur der inversen Laplace-Transformation zu verstehen, lohnt sich ein kurzer Blick auf ihren Ursprung: die Fourier-Transformation.
Die Fourier-Transformation zerlegt ein Signal in seine Frequenzanteile, setzt jedoch voraus, dass das Signal absolut integrierbar ist.
Viele physikalisch relevante Funktionen, wie etwa exponentiell wachsende Signale, erfüllen diese Bedingung nicht; ihr Fourier-Integral divergiert.

Um dieses Konvergenzproblem zu lösen, bedient man sich eines eleganten Tricks.
Man multipliziert die ursprüngliche Funktion $f(t)$ mit einem exponentiellen Dämpfungsfaktor $e^{-\gamma t}$, wobei $\gamma$ eine reelle Konstante ist.
Wählt man $\gamma$ gross genug, wird das Wachstum von $f(t)$ unterdrückt.
Wir definieren also eine gedämpfte Hilfsfunktion
\begin{equation}
    g(t) = f(t) e^{-\gamma t}.
    \label{laplace:eq:hilfsfunktion}
\end{equation}
Für diese gedämpfte Funktion existiert nun die Fourier-Transformation $G(\omega) = \mathcal{F}\{g(t)\}$:
\begin{equation}
    G(\omega) = \int_{-\infty}^{\infty} \left( f(t) e^{-\gamma t} \right) e^{-i\omega t} \, dt = \int_{-\infty}^{\infty} f(t) e^{-(\gamma + i\omega)t} \, dt.
    \label{laplace:eq:fourier_g}
\end{equation}
An dieser Stelle kommt uns die Physik zu Hilfe.
Die von uns betrachteten Vorgänge beginnen zu einem definierten Zeitpunkt, den wir als $t = 0$ wählen; für $t < 0$ setzen wir $f(t) = 0$ und nennen ein solches Signal \emph{kausal}.
Für kausale Funktionen verschwindet der Integrand in Gleichung~\eqref{laplace:eq:fourier_g} für alle negativen Zeiten, sodass sich das Integral effektiv nur noch von $0$ bis $\infty$ erstreckt.
Führen wir nun die komplexe Variable $s = \gamma + i\omega$ ein, so entspricht dieses Integral exakt der Definition der Laplace-Transformation von $f(t)$.
Es gilt also $G(\omega) = F(s)$.

Um nun $f(t)$ zurückzugewinnen, wenden wir die inverse Fourier-Transformation auf $G(\omega)$ an:
\begin{equation}
    g(t) = \frac{1}{2\pi} \int_{-\infty}^{\infty} G(\omega) e^{i\omega t} \, d\omega.
    \label{laplace:eq:inv_fourier_g}
\end{equation}
Setzen wir in diese Gleichung $g(t) = f(t) e^{-\gamma t}$ wieder ein und multiplizieren beide Seiten mit $e^{\gamma t}$, erhalten wir
\begin{equation}
    f(t) = \frac{e^{\gamma t}}{2\pi} \int_{-\infty}^{\infty} G(\omega) e^{i\omega t} \, d\omega 
         = \frac{1}{2\pi} \int_{-\infty}^{\infty} G(\omega) e^{(\gamma + i\omega)t} \, d\omega.
    \label{laplace:eq:rueck_1}
\end{equation}
Im letzten Schritt ersetzen wir die Integrationsvariable $\omega$ wieder durch unsere komplexe Variable $s$.
Mit $s = \gamma + i\omega$ und konstantem $\gamma$ folgt für das Differential $ds = i \, d\omega$, also $d\omega = \frac{ds}{i}$.
Da $\omega$ von $-\infty$ bis $\infty$ läuft, bewegt sich $s$ entlang einer vertikalen Geraden in der komplexen Ebene von $\gamma - i\infty$ bis $\gamma + i\infty$.
Setzen wir dies sowie $G(\omega) = F(s)$ in Gleichung~\eqref{laplace:eq:rueck_1} ein, so erhalten wir die fundamentale Formel für die inverse Laplace-Transformation:
\begin{equation}
    f(t) = \frac{1}{2\pi i} \int_{\gamma - i\infty}^{\gamma + i\infty} F(s) e^{st} \, ds.
    \label{laplace:eq:bromwich_integral}
\end{equation}

Dieses komplexe Linienintegral wird als \emph{Bromwich-Integral} bezeichnet.
Abbildung~\ref{laplace:fig:bromwich} veranschaulicht die geometrische Situation in der komplexen Ebene.
Die vertikale Integrationsgerade $\text{Re}(s) = \gamma$ wird dabei stets so gewählt, dass sie rechts von allen Singularitäten der Bildfunktion $F(s)$ liegt.
Man nennt den minimalen Wert, für den das Integral noch konvergiert, die Konvergenzabszisse $\gamma_0$.
Für alle $\gamma > \gamma_0$ liefert das Integral die eindeutige Lösung.

\begin{figure}
\centering
\scalebox{0.85}{
\begin{tikzpicture}

% ==============================
% 1. FARBEN DEFINIEREN
% ==============================
\definecolor{plotblue}{RGB}{0, 0, 200}
\definecolor{plotred}{RGB}{220, 30, 30}
\definecolor{plotgray}{RGB}{140, 140, 140}
\definecolor{plotbg}{RGB}{244, 244, 252}

% ==============================
% 2. HINTERGRUND (Konvergenzbereich)
% ==============================
\fill[plotbg] (1.5, -4.8) rectangle (5.0, 4.8);

% ==============================
% 3. ACHSEN
% ==============================
\draw[thick, -{Stealth[length=3.5mm, width=2.5mm]}] (-5.2, 0) -- (5.5, 0) node[right] {\textbf{Re}$(s)$};
\draw[thick, -{Stealth[length=3.5mm, width=2.5mm]}] (0, -5.2) -- (0, 5.2) node[above] {\textbf{Im}$(s)$};

% ==============================
% 4. LINIEN & BESCHRIFTUNGEN
% ==============================
% Konvergenzabszisse (\gamma_0)
\draw[thick, plotgray, dashed] (1.5, -4.8) -- (1.5, 4.8);
\draw[thick] (1.5, 0.15) -- (1.5, -0.15) node[below right, inner sep=2pt] {$\boldsymbol{\gamma_0}$};

% Bromwich-Pfad (\gamma)
\draw[thick, plotblue, ->-=0.55] (2.5, -4.8) -- (2.5, 4.8);
\draw[thick] (2.5, 0.15) -- (2.5, -0.15) node[below right, inner sep=2pt] {$\boldsymbol{\gamma}$};

% Text im Konvergenzbereich
\node[plotblue, align=center, font=\small] at (3.75, 3) {Konvergenz-\\bereich};

% ==============================
% 5. POLE (Rote Kreuze)
% ==============================
\newcommand{\pole}[2]{
    \draw[thick, plotred, line cap=round] (#1-0.12,#2-0.12) -- (#1+0.12,#2+0.12);
    \draw[thick, plotred, line cap=round] (#1-0.12,#2+0.12) -- (#1+0.12,#2-0.12);
}

% Pole platzieren
\pole{1.5}{0}
\pole{0}{0}
\pole{0}{2}
\pole{0}{-2}
\pole{-1.2}{1.5}
\pole{-1.2}{-1.5}
\pole{-2}{0}
\pole{-3}{0}

% ==============================
% 6. LEGENDE
% ==============================
\node[draw, thick, rounded corners=3pt, fill=white, inner xsep=5pt, inner ysep=4pt, anchor=north west] at (-5.0, 4.8) {
    \renewcommand{\arraystretch}{1.2}
    \begin{tabular}{@{}c l@{}}
        
        % Symbol 1: Polstelle
        \tikz[baseline=-0.6ex]{
            \draw[thick, plotred, line cap=round] (-0.1,-0.1) -- (0.1,0.1) (-0.1,0.1) -- (0.1,-0.1);
        } & Polstelle \\
        
        % Symbol 2: Konvergenzabszisse (etwas verlängert auf -0.3 bis 0.3)
        \tikz[baseline=-0.6ex]{
            \draw[thick, plotgray, dashed] (-0.3, 0) -- (0.3, 0);
        } & Konvergenzabszisse \\
        
        % Symbol 3: Bromwich-Pfad (Mit Dekoration für perfekten Strich)
        \tikz[baseline=-0.6ex]{
            \draw[thick, plotblue, ->-=0.6] (-0.3, 0) -- (0.3, 0); % Durchgehender Strich
        } & Bromwich-Pfad \\
        
    \end{tabular}
};

\end{tikzpicture}
}
\caption{Die Bromwich-Kontur in der komplexen $s$-Ebene. Die Integrationsgerade (blau) liegt rechts von der Konvergenzabszisse $\gamma_0$ und allen Singularitäten (rote Kreuze).}
\label{laplace:fig:bromwich}
\end{figure}


%
% Teil 2: Auswertung des Bromwich-Integrals mit dem Residuensatz
% und dem Lemma von Jordan (Schliessen der Kontur ueber den im Ursprung
% zentrierten Kreisbogen, D-Kontur)
%
\section{Auswertung mit dem Residuensatz: die geschlossene D-Kontur}
\label{laplace:subsection:residuensatz}

Mit Gleichung~\eqref{laplace:eq:bromwich_integral} ist die Inversion zwar
formal gelöst, praktisch stehen wir aber vor einem unendlich langen
Integrationsweg mitten in der komplexen Ebene.
Der Schlüssel zur konkreten Auswertung ist der Residuensatz, der im ersten
Teil dieses Buches ausführlich hergeleitet wurde und dessen Anwendung auf
Konturintegrale beispielsweise in \cite{laplace:ablowitz2003} vertieft wird.
Ist $G(s)$ innerhalb einer geschlossenen, positiv orientierten Kurve $C$
holomorph bis auf endlich viele isolierte Singularitäten $s_1,\dots,s_n$,
so gilt
\begin{equation}
    \oint_C G(s) \, ds
    =
    2\pi i \sum_{k=1}^{n} \operatorname{Res}\bigl(G(s), s_k\bigr).
    \label{laplace:eq:residuensatz}
\end{equation}
Entscheidend ist dabei ein oft übersehenes Detail:
Nur die von der Kurve \emph{eingeschlossenen} Singularitäten tragen zum
Integral bei.
Eine Polstelle, die auch nur knapp ausserhalb der Kontur liegt, ist für das
Integral schlicht unsichtbar.
Abbildung~\ref{laplace:fig:pol_auswahl} veranschaulicht diese Auswahlregel.

\begin{figure}
    \centering
    \scalebox{0.85}{
\begin{tikzpicture}

% ==============================
% 1. FARBEN DEFINIEREN
% ==============================
\definecolor{plotblue}{RGB}{0, 0, 200}
\definecolor{plotred}{RGB}{220, 30, 30}
\definecolor{plotgray}{RGB}{140, 140, 140}
\definecolor{plotbg}{RGB}{244, 244, 252}

% ==============================
% 2. ACHSEN
% ==============================
\draw[thick, -{Stealth[length=3.5mm, width=2.5mm]}] (-4.2, 0) -- (4.5, 0) node[right] {\textbf{Re}$(s)$};
\draw[thick, -{Stealth[length=3.5mm, width=2.5mm]}] (0, -4.2) -- (0, 4.2) node[above] {\textbf{Im}$(s)$};

% ==============================
% 3. KONTUR C
% ==============================
% Geschlossene, organische Blob-Form, positiv orientiert
\draw[thick, plotblue, decoration={markings, mark=at position 0.25 with {\arrow{Stealth[length=3mm, width=2mm]}},
                          mark=at position 0.60 with {\arrow{Stealth[length=3mm, width=2mm]}},
                          mark=at position 0.90 with {\arrow{Stealth[length=3mm, width=2mm]}}}, postaction={decorate}] 
    (-0.5, 0) .. controls (1, 3) and (3, 1) .. 
    (2, -1) .. controls (1.5, -3) and (-3, -3) .. 
    (-2.5, -1) .. controls (-2, 1) and (-1, -1.5) .. (-0.5, 0);

% Label C aussen
\node[plotblue, font=\large] at (1.8, 1.8) {$C$};

% ==============================
% 4. POLE (Kreuze)
% ==============================
\newcommand{\pole}[3][plotred]{
    \draw[thick, #1, line cap=round] (#2-0.12,#3-0.12) -- (#2+0.12,#3+0.12);
    \draw[thick, #1, line cap=round] (#2-0.12,#3+0.12) -- (#2+0.12,#3-0.12);
}

% Innere Pole (rot)
\pole{0.8}{0.8}
\pole{0.0}{-1.8}
\pole{-1.0}{-1.2}

% Aeussere Pole (grau)
\pole[plotgray]{3.0}{2.0}
\pole[plotgray]{-2.8}{-2.0}

% ==============================
% 5. LEGENDE
% ==============================
\node[draw, thick, rounded corners=3pt, fill=white, inner xsep=5pt, inner ysep=4pt, anchor=north west] at (-4.0, 4.0) {
    \renewcommand{\arraystretch}{1.2}
    \begin{tabular}{@{}c l@{}}
        
        \tikz[baseline=-0.6ex]{
            \draw[thick, plotred, line cap=round] (-0.1,-0.1) -- (0.1,0.1) (-0.1,0.1) -- (0.1,-0.1);
        } & Relevant \\
        
        \tikz[baseline=-0.6ex]{
            \draw[thick, plotgray, line cap=round] (-0.1,-0.1) -- (0.1,0.1) (-0.1,0.1) -- (0.1,-0.1);
        } & Irrelevant \\
        
        \tikz[baseline=-0.6ex]{
            \draw[thick, plotblue, ->-=0.6] (-0.3, 0) -- (0.3, 0);
        } & Kontur $C$ \\
        
    \end{tabular}
};

\end{tikzpicture}
}

    \caption{Auswahlregel des Residuensatzes:
    Nur die von der geschlossenen Kontur $C$ eingeschlossenen Polstellen
    (rot) liefern einen Beitrag $2\pi i \operatorname{Res}$ zum
    Konturintegral.
    Die Polstelle ausserhalb der Kontur (grau) ist für das Integral
    unsichtbar.}
    \label{laplace:fig:pol_auswahl}
\end{figure}

Der Residuensatz verlangt allerdings eine \emph{geschlossene} Kurve, während
der Bromwich-Pfad eine offene, unendlich lange Gerade ist.
Wir müssen die Gerade also künstlich schliessen, und zwar so, dass der
hinzugefügte Wegabschnitt im Grenzfall nichts zum Integral beiträgt.
Die Richtung des Schliessens gibt der Faktor $e^{st}$ vor.
Wegen
\begin{equation}
    \bigl| e^{st} \bigr|
    =
    e^{t \, \text{Re}(s)}
    \label{laplace:eq:betrag_est}
\end{equation}
fällt der Integrand für $t > 0$ exponentiell ab, sobald wir uns in die linke
Halbebene $\text{Re}(s) < 0$ begeben, und er explodiert in der rechten.
Wir schliessen die Kontur daher nach links, und zwar über einen grossen,
im Ursprung zentrierten Kreisbogen $\Gamma_R$, wie ihn auch
\cite{laplace:carslaw1959} verwendet.
Da sämtliche Singularitäten von $F(s)$ links der Geraden
$\text{Re}(s) = \gamma$ liegen, fängt die geschlossene Kurve sie für
hinreichend grosses $R$ alle ein.
Zusammen mit dem vertikalen Geradenstück, das als Sehne des Kreises
zwischen den beiden Schnittpunkten mit dem Bogen verläuft, entsteht die
in Abbildung~\ref{laplace:fig:d_kontur} gezeigte geschlossene Kurve, die
wegen ihrer Form als \emph{D-Kontur} bezeichnet wird.
Dass der ursprungszentrierte Bogen dabei im Streifen
$0 < \text{Re}(s) < \gamma$ ein kleines Stück in die rechte Halbebene
hineinragt, wird sich gleich als harmlos erweisen.

\begin{figure}
    \centering
    \scalebox{0.85}{
\begin{tikzpicture}

% ==============================
% 1. FARBEN DEFINIEREN
% ==============================
\definecolor{plotblue}{RGB}{0, 0, 200}
\definecolor{plotred}{RGB}{220, 30, 30}
\definecolor{plotgray}{RGB}{140, 140, 140}
\definecolor{plotorange}{RGB}{240, 140, 0}
\definecolor{plotbg}{RGB}{244, 244, 252}

% ==============================
% 2. HINTERGRUND FUELLEN
% ==============================
%\fill[plotbg] (1.5, -3.923) -- (1.5, 3.923) arc (69.07:290.93:4.2) -- cycle;

% ==============================
% 3. ACHSEN
% ==============================
\draw[thick, -{Stealth[length=3.5mm, width=2.5mm]}] (-5.5, 0) -- (3.5, 0) node[right] {\textbf{Re}$(s)$};
\draw[thick, -{Stealth[length=3.5mm, width=2.5mm]}] (0, -5.0) -- (0, 5.5) node[above] {\textbf{Im}$(s)$};

% ==============================
% 4. D-KONTUR
% ==============================

% Bromwich-Pfad (Sehne)
\draw[thick, plotblue, decoration={markings, mark=at position 0.55 with {\arrow{Stealth[length=3mm, width=2mm]}}}, postaction={decorate}] 
    (1.5, -3.923) -- (1.5, 3.923);

% Gamma Tick
\draw[thick] (1.5, 0.1) -- (1.5, -0.1) node[below right, inner sep=2pt] {$\boldsymbol{\gamma}$};

% Ticks an den Enden des Bromwich-Pfades
% (Schnittpunkte Sehne/Kreis pythagoreisch exakt: gamma +- i*sqrt(R^2-gamma^2);
%  fuer R -> infty naehern sie sich gamma +- iR an)
\draw[thick] (1.35, 3.923) -- (1.65, 3.923) node[right, text=black] {$\gamma + i\sqrt{R^2 - \gamma^2}$};
\draw[thick] (1.35, -3.923) -- (1.65, -3.923) node[right, text=black] {$\gamma - i\sqrt{R^2 - \gamma^2}$};

% Kreisbogen Gamma_R
\draw[thick, plotorange, decoration={markings, 
        mark=at position 0.3 with {\arrow{Stealth[length=3mm, width=2mm]}},
        mark=at position 0.7 with {\arrow{Stealth[length=3mm, width=2mm]}}}, 
        postaction={decorate}] 
    (1.5, 3.923) arc (69.07:290.93:4.2);

\node[plotorange, font=\large] at (-4.6, 1.8) {$\Gamma_R$};

% Markierung fuer den Bereich in der rechten Halbebene
% Ragt in rechte Halbebene (entfernt)

% Radiuspfeil
\draw[thick, plotgray, dashed, -{Stealth[length=2.5mm, width=2mm]}] (0,0) -- (-2.97, 2.97) node[midway, below left] {$R\to\infty$};

% ==============================
% 5. POLE (Kreuze)
% ==============================
\newcommand{\pole}[3][plotred]{
    \draw[thick, #1, line cap=round] (#2-0.12,#3-0.12) -- (#2+0.12,#3+0.12);
    \draw[thick, #1, line cap=round] (#2-0.12,#3+0.12) -- (#2+0.12,#3-0.12);
}

% Innere Pole (rot)
\pole{0.0}{2.0}
\pole{0.0}{-2.0}
\pole{-1.2}{0.9}
\pole{-2.0}{-0.5}
\pole{-0.5}{0.0}

% ==============================
% 6. LEGENDE
% ==============================
\node[draw, thick, rounded corners=3pt, fill=white, inner xsep=5pt, inner ysep=4pt, anchor=north west] at (-5.3, 5.3) {
    \renewcommand{\arraystretch}{1.2}
    \begin{tabular}{@{}c l@{}}
        
        \tikz[baseline=-0.6ex]{
            \draw[thick, plotred, line cap=round] (-0.1,-0.1) -- (0.1,0.1) (-0.1,0.1) -- (0.1,-0.1);
        } & Polstelle \\
        
        \tikz[baseline=-0.6ex]{
            \draw[thick, plotblue, decoration={markings, mark=at position 0.55 with {\arrow{Stealth[length=2mm, width=1.5mm]}}}, postaction={decorate}] (-0.3, 0) -- (0.3, 0);
        } & Bromwich-Pfad \\
        
        \tikz[baseline=-0.6ex]{
            \draw[thick, plotorange, decoration={markings, mark=at position 0.55 with {\arrow{Stealth[length=2mm, width=1.5mm]}}}, postaction={decorate}] (-0.3, 0) -- (0.3, 0);
        } & Kreisbogen $\Gamma_R$ \\
        
    \end{tabular}
};

\end{tikzpicture}
}

    \caption{Die geschlossene D-Kontur:
    Der vertikale Bromwich-Pfad bildet die Sehne eines im Ursprung
    zentrierten Kreises, dessen grosser Bogen $\Gamma_R$ die Kurve nach
    links schliesst und dabei im Streifen $0 < \text{Re}(s) < \gamma$
    leicht in die rechte Halbebene hineinragt.
    Für $R \to \infty$ umschliesst die Kurve alle Singularitäten von
    $F(s)$, während der Beitrag des Bogens nach dem Lemma von Jordan
    verschwindet.}
    \label{laplace:fig:d_kontur}
\end{figure}

Dass der Bogenbeitrag tatsächlich verschwindet, ist keineswegs
selbstverständlich, denn die Bogenlänge wächst proportional zu $R$ ins
Unendliche.
Tief in der linken Halbebene ist das unkritisch:
Dort erdrückt der exponentielle Abfall von $e^{st}$ aus
Gleichung~\eqref{laplace:eq:betrag_est} das nur lineare Anwachsen der
Weglänge mühelos.
Kritisch sind allein die Bogenstücke nahe der imaginären Achse und im
schmalen Streifen bis $\text{Re}(s) = \gamma$, wo
$|e^{st}| \leq e^{\gamma t}$ lediglich beschränkt ist; dort muss der
gleichmässige Abfall der Bildfunktion einspringen.
Genau das leistet das \emph{Lemma von Jordan} in der für die
Laplace-Inversion passenden Form, wie sie etwa in
\cite{laplace:doetsch1961} als Satz 34.2 bewiesen wird:
Strebt $F(s)$ auf den Bögen für $|s| \to \infty$ gleichmässig gegen
null, so gilt für $t > 0$
\begin{equation}
    \lim_{R \to \infty} \int_{\Gamma_R} e^{st} F(s) \, ds
    =
    0.
    \label{laplace:eq:jordan}
\end{equation}
Das Resultat deckt auch die rechts der imaginären Achse liegenden
Bogenstücke ab, denn deren geometrische Länge bleibt für $R \to \infty$
beschränkt:
Jedes der beiden Bogenstücke durchquert den Streifen
$0 < \text{Re}(s) < \gamma$ nahezu senkrecht, und seine Länge
$R \arcsin(\gamma/R)$ strebt für $R \to \infty$ gegen den festen Wert
$\gamma$.
Auf einem Wegstück beschränkter Länge, auf dem $|e^{st}| \leq e^{\gamma t}$
gilt, löscht der gleichmässige Abfall von $F(s)$ den Beitrag ebenfalls
aus.
Für die in der Praxis auftretenden Bildfunktionen, etwa gebrochen rationale
Funktionen mit Nennergrad grösser als Zählergrad, ist diese Voraussetzung
stets erfüllt.

Damit fällt das Puzzle zusammen.
Lassen wir in der D-Kontur den Radius $R$ gegen unendlich streben, so
verschwindet der Bogenbeitrag nach
Gleichung~\eqref{laplace:eq:jordan}, das vertikale Geradenstück wird zum
Bromwich-Integral, und der Residuensatz
\eqref{laplace:eq:residuensatz} liefert die bemerkenswert kompakte
Inversionsformel
\begin{equation}
    f(t)
    =
    \frac{1}{2\pi i} \int_{\gamma - i\infty}^{\gamma + i\infty}
    F(s) e^{st} \, ds
    =
    \sum_{k} \operatorname{Res}\bigl(F(s) e^{st}, s_k\bigr).
    \label{laplace:eq:inversion_residuen}
\end{equation}
Das unendliche Linienintegral ist zu einer endlichen Summe über die
Singularitäten von $F(s)$ geschrumpft.
Die gesamte Information über den zeitlichen Verlauf von $f(t)$ steckt also
in der Lage und Art der Singularitäten der Bildfunktion.

Ein kurzer Test an einem bekannten Korrespondenzpaar zeigt, dass die Formel
hält, was sie verspricht.
Die Bildfunktion $F(s) = \frac{1}{s-a}$ besitzt einen einfachen Pol bei
$s = a$ mit dem Residuum
\begin{equation}
    \operatorname{Res}\left(\frac{e^{st}}{s-a}, a\right)
    =
    \lim_{s \to a} (s-a) \frac{e^{st}}{s-a}
    =
    e^{at}.
    \label{laplace:eq:beispiel_exponential}
\end{equation}
Gleichung~\eqref{laplace:eq:inversion_residuen} liefert also
$f(t) = e^{at}$ und reproduziert damit exakt den Eintrag der
Exponentialfunktion aus Tabelle~\ref{laplace:tab:korrespondenzen}.
Was dort als auswendig gelernte Korrespondenz stand, ist jetzt das Residuum
eines Konturintegrals.

So elegant dieser Mechanismus ist, er trägt eine versteckte Voraussetzung
in sich:
$F(s)$ muss in der ganzen linken Halbebene eine \emph{eindeutige},
meromorphe Funktion sein, denn nur dann dürfen wir die D-Kontur überhaupt
zuziehen.
Genau diese Eindeutigkeit geht verloren, sobald Wurzeln oder Logarithmen
der Variablen $s$ auftreten.
Warum das bei physikalisch völlig harmlosen Problemen passiert und wie die
Kontur dann repariert werden muss, ist das Thema des nächsten Abschnitts.

%
% Teil 3: Verzweigungspunkte und Verzweigungsschnitte am Beispiel
% des halbunendlichen Waermeleiters; Schluesselloch-Kontur (Keyhole)
%
\section{Verzweigungsschnitte und die Schlüsselloch-Kontur}
\label{laplace:subsection:verzweigung}

Wie schnell mehrdeutige Funktionen auftreten, zeigt das Paradebeispiel der
instationären Wärmeleitung.
Ein halbunendlicher Stab $x \geq 0$ mit Anfangstemperatur null wird durch
die normierte Wärmeleitungsgleichung
\begin{equation}
    \frac{\partial^2 U}{\partial x^2}
    =
    \frac{\partial U}{\partial t}
    \label{laplace:eq:waermeleitung}
\end{equation}
beschrieben, wobei am Ende $x = 0$ eine vorgegebene Randtemperatur
$A_0(t)$ eingeprägt wird und die verschwindende Anfangstemperatur formal
durch die Anfangsbedingung
\begin{equation}
    U(x, 0)
    =
    0
    \label{laplace:eq:anfangsbedingung}
\end{equation}
ausgedrückt wird.
Die Laplace-Transformation bezüglich der Zeit macht aus der Zeitableitung
gemäss der Ableitungsregel aus
Tabelle~\ref{laplace:tab:korrespondenzen} eine Multiplikation,
\begin{equation}
    \mathcal{L}\left\{\frac{\partial U}{\partial t}\right\}
    =
    s \, u(x,s) - U(x,0)
    =
    s \, u(x,s),
    \label{laplace:eq:ableitungsregel}
\end{equation}
wobei der Anfangswertterm genau wegen
Gleichung~\eqref{laplace:eq:anfangsbedingung} entfällt.
Aus der partiellen Differentialgleichung wird so eine gewöhnliche
Differentialgleichung für die Bildfunktion
$u(x,s) = \mathcal{L}\{U(x,t)\}$, nämlich
\begin{equation}
    \frac{d^2 u}{d x^2}
    =
    s \, u,
    \label{laplace:eq:bildgleichung}
\end{equation}
mit der allgemeinen Lösung
$u(x,s) = c_1 e^{x\sqrt{s}} + c_2 e^{-x\sqrt{s}}$.
Da die Temperatur für $x \to \infty$ beschränkt bleiben muss, scheidet der
anwachsende Anteil aus; das setzt allerdings $\text{Re}\sqrt{s} > 0$
voraus, denn nur dann wächst $e^{x\sqrt{s}}$ tatsächlich an, während
$e^{-x\sqrt{s}}$ beschränkt bleibt.
Auf dem gleich eingeführten Hauptzweig der Wurzel ist genau das der Fall,
sodass $c_1 = 0$ folgt.
Mit der transformierten Randbedingung
$a_0(s) = \mathcal{L}\{A_0(t)\}$ erhalten wir
\begin{equation}
    u(x,s)
    =
    a_0(s) \, u_0(x,s),
    \qquad
    u_0(x,s)
    =
    e^{-x\sqrt{s}},
    \label{laplace:eq:halbunendlich_loesung}
\end{equation}
wie in \cite{laplace:doetsch1961} im Detail hergeleitet wird.
Die Lösung zerfällt also in die transformierte Randanregung $a_0(s)$ und
eine \emph{Basislösung} $u_0(x,s)$, die am Rand $x = 0$ den Wert eins
annimmt und das reine Ortsprofil beschreibt; dieselbe Aufspaltung wird uns
beim Stab endlicher Länge wieder begegnen.
Die Geometrie des halbunendlichen Gebietes hat uns also zwangsläufig eine
Wurzel $\sqrt{s}$ eingebrockt.

Diese Wurzel ist keine harmlose Funktion, sondern \emph{mehrdeutig}.
Zu jedem $s \neq 0$ existieren zwei Werte $w$ mit $w^2 = s$, und welcher
davon gemeint ist, hängt vom gewählten Zweig ab.
Schreiben wir $s = r e^{i\theta}$ und umlaufen den Ursprung einmal, indem
$\theta$ von $0$ bis $2\pi$ wächst, so geht
$\sqrt{s} = \sqrt{r}\, e^{i\theta/2}$ in
$\sqrt{r}\, e^{i\pi} = -\sqrt{r}$ über:
Die Funktion kehrt nicht zu ihrem Ausgangswert zurück, sondern wechselt
das Vorzeichen \cite{laplace:ablowitz2003}.
Der Ursprung $s = 0$ ist deshalb ein \emph{Verzweigungspunkt} von
$\sqrt{s}$.
Für den Residuensatz ist das fatal, denn er setzt eine eindeutige,
holomorphe Funktion voraus.
Der Ausweg besteht darin, geschlossene Umläufe um den Verzweigungspunkt von
vornherein zu verbieten:
Wir schlitzen die Ebene entlang eines \emph{Verzweigungsschnittes} auf, der
den Verzweigungspunkt $s = 0$ mit dem zweiten Verzweigungspunkt
$s = \infty$ verbindet, und arbeiten auf dem Hauptzweig
$-\pi < \arg s < \pi$.
Der Schnitt liegt dann auf der negativen reellen Achse, und diese Wahl ist
kein Zufall.
Auf dem Hauptzweig gilt nämlich $\text{Re}\sqrt{s} \geq 0$, sodass
$|e^{-x\sqrt{s}}| \leq 1$ beschränkt bleibt, was genau der physikalischen
Beschränktheit der Temperatur entspricht.
Ausserdem bleibt die rechte Halbebene mitsamt dem Bromwich-Pfad
unangetastet.

Nur dürfen wir die D-Kontur aus
Abschnitt~\ref{laplace:subsection:residuensatz} jetzt nicht mehr einfach
zuziehen, denn ihr Kreisbogen würde den Schnitt bei $s = -R$ einmal
durchstossen.
Die Kontur muss dem Schnitt ausweichen:
Sie läuft von oben kommend bis knapp an die negative reelle Achse heran,
folgt ihr oberhalb des Schnittes nach innen, umrundet den
Verzweigungspunkt auf einem kleinen Kreis $\Gamma_\rho$ und kehrt
unterhalb des Schnittes wieder nach aussen zurück.
Das Ergebnis ist die berühmte \emph{Schlüsselloch-Kontur}, die in
Abbildung~\ref{laplace:fig:keyhole_kontur} dargestellt ist.
Innerhalb dieser Kurve ist der Integrand eindeutig und holomorph bis auf
isolierte Pole, sodass der Residuensatz wieder anwendbar ist.

\begin{figure}
    \centering
    \scalebox{0.8}{
\begin{tikzpicture}

% ==============================
% 1. FARBEN DEFINIEREN
% ==============================
\definecolor{plotblue}{RGB}{0, 0, 200}
\definecolor{plotred}{RGB}{220, 30, 30}
\definecolor{plotgray}{RGB}{140, 140, 140}
\definecolor{plotbg}{RGB}{244, 244, 252}
\definecolor{plotorange}{RGB}{230, 140, 0}
\definecolor{plotgreen}{RGB}{0, 150, 80}
\definecolor{plotpurple}{RGB}{140, 60, 180}

% ==============================
% 2. ACHSEN UND SCHNITT
% ==============================
% Verzweigungsschnitt (Wellenlinie)
\draw[thick, plotred, decorate, decoration={snake, amplitude=0.8mm, segment length=3mm}] (-5.1, 0) -- (0, 0);

\draw[thick, -{Stealth[length=3.5mm, width=2.5mm]}] (-5.5, 0) -- (2.5, 0) node[right] {\textbf{Re}$(s)$};
\draw[thick, -{Stealth[length=3.5mm, width=2.5mm]}] (0, -5.0) -- (0, 5.0) node[above] {\textbf{Im}$(s)$};

% Verzweigungspunkt
\fill[plotred] (0,0) circle (2.5pt);

% ==============================
% 3. KONTUR-SEGMENTE
% ==============================
% 1. Bromwich-Pfad (Blau)
\draw[thick, plotblue, decoration={markings, mark=at position 0.6 with {\arrow{Stealth[length=3mm, width=2mm]}}}, postaction={decorate}] 
    (1.2, -4.025) -- (1.2, 4.025);
\draw[thick] (1.2, 0.1) -- (1.2, -0.1) node[below right, inner sep=2pt] {$\boldsymbol{\gamma}$};

% 2. Gamma_R oben (Orange)
\draw[thick, plotorange, decoration={markings, mark=at position 0.5 with {\arrow{Stealth[length=3mm, width=2mm]}}}, postaction={decorate}] 
    (1.2, 4.025) arc (73.39:177.95:4.2);
\node[plotorange, font=\large] at (-3.5, 3.5) {$\Gamma_R$};

% 3. L_+ (Gruen)
\draw[thick, plotgreen, decoration={markings, mark=at position 0.5 with {\arrow{Stealth[length=3mm, width=2mm]}}}, postaction={decorate}] 
    (-4.197, 0.15) -- (-0.37, 0.15);
\node[plotgreen, anchor=south] at (-2.0, 0.65) {$L_+$};

% 4. Gamma_rho (Violett)
\draw[thick, plotpurple, decoration={markings, mark=at position 0.5 with {\arrow{Stealth[length=3mm, width=2mm]}}}, postaction={decorate}] 
    (-0.37, 0.15) arc (157.98:-157.98:0.4);
\node[plotpurple, font=\large] at (0.7, 0.5) {$\Gamma_\rho$};

% 5. L_- (Gruen)
\draw[thick, plotgreen, decoration={markings, mark=at position 0.5 with {\arrow{Stealth[length=3mm, width=2mm]}}}, postaction={decorate}] 
    (-0.37, -0.15) -- (-4.197, -0.15);
\node[plotgreen, anchor=north] at (-2.0, -0.65) {$L_-$};

% 6. Gamma_R unten (Orange)
\draw[thick, plotorange, decoration={markings, mark=at position 0.5 with {\arrow{Stealth[length=3mm, width=2mm]}}}, postaction={decorate}] 
    (-4.197, -0.15) arc (182.05:286.61:4.2);

% ==============================
% 4. ZWEIG-BESCHRIFTUNGEN
% ==============================
\node[plotgray, font=\small, anchor=south] at (-2.0, 0.2) {$s=re^{i\pi},\ \sqrt{s}=i\sqrt{r}$};
\node[plotgray, font=\small, anchor=north] at (-2.0, -0.2) {$s=re^{-i\pi},\ \sqrt{s}=-i\sqrt{r}$};

% ==============================
% 5. POLE (Kreuze)
% ==============================
\newcommand{\pole}[3][plotred]{
    \draw[thick, #1, line cap=round] (#2-0.15,#3-0.15) -- (#2+0.15,#3+0.15);
    \draw[thick, #1, line cap=round] (#2-0.15,#3+0.15) -- (#2+0.15,#3-0.15);
}

% Pole auf imag. Achse
\pole{0.0}{2.2}
\node[plotred, right] at (0.15, 2.2) {$i\omega$};

\pole{0.0}{-2.2}
\node[plotred, right] at (0.15, -2.2) {$-i\omega$};

% ==============================
% 6. LEGENDE
% ==============================
\node[draw, thick, rounded corners=3pt, fill=white, inner xsep=5pt, inner ysep=4pt, anchor=west] at (3.8, 0) {
    \renewcommand{\arraystretch}{1.2}
    \begin{tabular}{@{}c l@{}}
        
        \tikz[baseline=-0.6ex]{
            \fill[plotred] (0,0) circle (2pt);
        } & Verzweigungspunkt \\
        
        \tikz[baseline=-0.6ex]{
            \draw[thick, plotred] (-0.45, 0) .. controls (-0.35, 0.2) and (-0.25, -0.2) .. (-0.15, 0) .. controls (-0.05, 0.2) and (0.05, -0.2) .. (0.15, 0) .. controls (0.25, 0.2) and (0.35, -0.2) .. (0.45, 0);
        } & Verzweigungsschnitt \\
        
        \tikz[baseline=-0.6ex]{
            \draw[thick, plotred, line cap=round] (-0.1,-0.1) -- (0.1,0.1) (-0.1,0.1) -- (0.1,-0.1);
        } & Polstelle \\
        
        \tikz[baseline=-0.6ex]{
            \draw[thick, plotblue, ->-=0.55] (-0.3,0) -- (0.3,0);
        } & Bromwich-Pfad \\
        
        \tikz[baseline=-0.6ex]{
            \draw[thick, plotorange, ->-=0.55] (-0.3,0) -- (0.3,0);
        } & Kreisbogen $\Gamma_R$ \\
        
        \tikz[baseline=-0.6ex]{
            \draw[thick, plotgreen, ->-=0.55] (-0.3,0) -- (0.3,0);
        } & Ufer $L_+, L_-$ \\
        
        \tikz[baseline=-0.6ex]{
            \draw[thick, plotpurple, ->-=0.55] (-0.3,0) -- (0.3,0);
        } & Kreisbogen $\Gamma_\rho$ \\
        
    \end{tabular}
};

\end{tikzpicture}
}

    \caption{Die Schlüsselloch-Kontur für Bildfunktionen mit
    Verzweigungspunkt bei $s = 0$.
    Der Verzweigungsschnitt (rote Wellenlinie) auf der negativen reellen
    Achse wird von den beiden horizontalen Wegen $L_+$ und $L_-$
    eingerahmt, der kleine Kreis $\Gamma_\rho$ umgeht den
    Verzweigungspunkt.
    Innerhalb der Kontur ist $\sqrt{s}$ eindeutig, und der Residuensatz
    erfasst nur die eingeschlossenen Pole, hier $s = \pm i\omega$.}
    \label{laplace:fig:keyhole_kontur}
\end{figure}

Wie diese Kontur konkret arbeitet, zeigt das klassische Beispiel aus
\cite{laplace:doetsch1961}, bei dem das Stabende mit der Temperatur
$A_0(t) = \cos \omega t$ angeregt wird.
Wegen $\mathcal{L}\{\cos \omega t\} = \frac{s}{s^2 + \omega^2}$ lautet die
Bildfunktion
\begin{equation}
    u(x,s)
    =
    \frac{s}{s^2 + \omega^2} \, e^{-x\sqrt{s}},
    \label{laplace:eq:keyhole_beispiel}
\end{equation}
mit einfachen Polen bei $s = \pm i\omega$ und dem Verzweigungspunkt bei
$s = 0$.
Der Residuensatz auf der Schlüsselloch-Kontur liefert $2\pi i$ mal die
Residuen an den beiden eingeschlossenen Polen.
Im Grenzübergang verschwinden die grossen Bögen nach dem Lemma von Jordan,
denn $\frac{s}{s^2+\omega^2}$ fällt wie $\frac{1}{s}$ ab und
$|e^{-x\sqrt{s}}| \leq 1$ stört dabei nicht.
Auch der kleine Kreis $\Gamma_\rho$ trägt nichts bei, weil der Integrand
dort beschränkt ist und die Weglänge gegen null schrumpft.
Die Residuen selbst rechnen wir in zwei Schritten aus.
Am einfachen Pol $s = i\omega$ liefert die Grenzwertformel zunächst
\begin{equation}
    \operatorname{Res}\bigl(u(x,s)e^{st}, i\omega\bigr)
    =
    \lim_{s \to i\omega}
    (s - i\omega) \,
    \frac{s \, e^{st - x\sqrt{s}}}{(s - i\omega)(s + i\omega)}
    =
    \frac{i\omega \, e^{i\omega t - x\sqrt{i\omega}}}{2 i\omega}
    =
    \frac{1}{2} \, e^{i\omega t - x\sqrt{i\omega}}.
    \label{laplace:eq:keyhole_grenzwert}
\end{equation}
Die verbleibende Wurzel werten wir in Polarform aus.
Wegen $i = e^{i\pi/2}$ ist
$\sqrt{i\omega}
= \sqrt{\omega} \, e^{i\pi/4}
= \sqrt{\omega} \, \frac{1 + i}{\sqrt{2}}
= \sqrt{\tfrac{\omega}{2}} \, (1 + i)$,
und am Pol $s = -i\omega$ ergibt sich völlig analog der konjugierte
Wert.
Zusammengefasst lauten die beiden Residuen also
\begin{equation}
    \operatorname{Res}\bigl(u(x,s)e^{st}, \pm i\omega\bigr)
    =
    \frac{1}{2} \,
    e^{-x\sqrt{\omega/2}} \,
    e^{\pm i\left(\omega t - x\sqrt{\omega/2}\right)},
    \label{laplace:eq:keyhole_residuen}
\end{equation}
und ihre Summe ergibt den reellen Ausdruck
$e^{-x\sqrt{\omega/2}} \cos\bigl(\omega t - x\sqrt{\omega/2}\bigr)$.

Der eigentliche Clou passiert aber am Schnitt.
Oberhalb gilt $s = r e^{i\pi}$ und damit $\sqrt{s} = i\sqrt{r}$, unterhalb
dagegen $s = r e^{-i\pi}$ und $\sqrt{s} = -i\sqrt{r}$.
Die Integranden auf $L_+$ und $L_-$ sind also verschieden, die beiden
Wegintegrale heben sich \emph{nicht} auf, und genau ihre Differenz bleibt
als reelles Integral übrig.
Sortiert man alle Beiträge, so entsteht die Lösung
\begin{equation}
    U(x,t)
    =
    e^{-x\sqrt{\omega/2}}
    \cos\Bigl(\omega t - x\sqrt{\tfrac{\omega}{2}}\Bigr)
    -
    \frac{1}{\pi}
    \int_0^\infty
    e^{-rt} \, \frac{r}{r^2 + \omega^2} \, \sin\bigl(x \sqrt{r}\bigr)
    \, dr.
    \label{laplace:eq:keyhole_loesung}
\end{equation}
Diese Formel erzählt die ganze Physik des Problems.
Der erste Term stammt von den Polen und beschreibt den eingeschwungenen
Zustand:
Eine Temperaturwelle dringt in den Stab ein, wird dabei exponentiell
gedämpft und läuft der Anregung mit wachsender Tiefe zunehmend hinterher.
Der zweite Term stammt vom Verzweigungsschnitt und ist ein transienter
Einschwingvorgang, der wegen des Faktors $e^{-rt}$ für grosse Zeiten
abklingt.
Während isolierte Pole diskrete Beiträge $e^{s_k t}$ liefern, wirkt der
Schnitt wie ein \emph{Kontinuum} von Abklingraten $e^{-rt}$, über das
integriert werden muss.

Damit drängt sich eine Frage auf.
Der Verzweigungsschnitt kam über die Wurzel $\sqrt{s}$ ins Spiel, und die
Wurzel kam über das halbunendliche Gebiet.
Was passiert, wenn der Stab eine endliche Länge $l$ besitzt?
Die Antwort ist verblüffend und der schönste Aha-Moment dieser Arbeit.

%
% Teil 4: Der Stab endlicher Laenge -- keine Verzweigung, sondern
% unendlich viele Pole auf der negativen reellen Achse;
% Residuenreihe (Fourier-Reihe) und Grenzuebergang l -> unendlich
%
\section{Der Stab endlicher Länge: unendlich viele Pole statt
Verzweigungsschnitt}
\label{laplace:subsection:endlicher_stab}

Wir betrachten nun denselben Wärmeleiter, aber mit endlicher Länge $l$.
Am linken Ende $x = 0$ wird die Temperatur $A_0(t)$ eingeprägt, das rechte
Ende $x = l$ wird auf null gehalten, und die Anfangstemperatur verschwindet.
Die Bildgleichung~\eqref{laplace:eq:bildgleichung} bleibt unverändert, nur
treten an die Stelle der Beschränktheitsforderung jetzt zwei
Randbedingungen.
Wie beim halbunendlichen Stab spalten wir die Lösung in die Randanregung
$a_0(s)$ und eine Basislösung $u_0$ auf, die am linken Rand den Wert eins
und am rechten den Wert null annimmt.
Die Basislösung gewinnen wir aus der allgemeinen exponentiellen Lösung
$c_1 e^{x\sqrt{s}} + c_2 e^{-x\sqrt{s}}$ von
Gleichung~\eqref{laplace:eq:bildgleichung}:
Die Randbedingung bei $x = l$ erzwingt die antisymmetrische Kombination
$\sinh\bigl((l-x)\sqrt{s}\bigr)
= \tfrac{1}{2}\bigl(e^{(l-x)\sqrt{s}} - e^{-(l-x)\sqrt{s}}\bigr)$,
die dort verschwindet, und die Normierung auf den Wert eins bei $x = 0$
liefert den Nenner $\sinh\bigl(l\sqrt{s}\bigr)$.
So ergibt sich
\begin{equation}
    u(x,s)
    =
    a_0(s) \, u_0(x,s),
    \qquad
    u_0(x,s)
    =
    \frac{\sinh\bigl((l-x)\sqrt{s}\bigr)}{\sinh\bigl(l\sqrt{s}\bigr)},
    \label{laplace:eq:stab_bildfunktion}
\end{equation}
wie ebenfalls in \cite{laplace:doetsch1961} vorgerechnet wird.
Um die Formeln übersichtlich zu halten, setzen wir im Folgenden $l = \pi$;
der allgemeine Fall entsteht daraus durch einfache Skalierung.

Wieder steht $\sqrt{s}$ in der Bildfunktion, und nach dem letzten Abschnitt
erwarten wir reflexartig einen Verzweigungsschnitt.
Doch die Erwartung trügt.
Ersetzen wir probeweise $\sqrt{s}$ durch $-\sqrt{s}$, was genau dem
Wechsel auf den anderen Zweig der Wurzel entspricht, so wechseln wegen
$\sinh(-z) = -\sinh(z)$ Zähler und Nenner \emph{gleichzeitig} das
Vorzeichen:
\begin{equation}
    \frac{\sinh\bigl(-(\pi-x)\sqrt{s}\bigr)}
         {\sinh\bigl(-\pi\sqrt{s}\bigr)}
    =
    \frac{-\sinh\bigl((\pi-x)\sqrt{s}\bigr)}
         {-\sinh\bigl(\pi\sqrt{s}\bigr)}
    =
    u_0(x,s).
    \label{laplace:eq:eindeutigkeit}
\end{equation}
Die Funktion $u_0$ merkt vom Zweigwechsel überhaupt nichts, sie ist
\emph{eindeutig}.
Noch deutlicher wird das in der Potenzreihe, denn $\sinh$ ist ungerade, und
deshalb enthält
\begin{equation}
    \sinh\bigl(a\sqrt{s}\bigr)
    =
    \sqrt{s}
    \left(
        a + \frac{a^3}{3!}\,s + \frac{a^5}{5!}\,s^2 + \cdots
    \right)
    \label{laplace:eq:sinh_reihe}
\end{equation}
die Wurzel nur als gemeinsamen Vorfaktor, der sich im Quotienten
$u_0(x,s)$ herauskürzt.
Übrig bleibt eine Potenzreihe in $s$ selbst.
Das ist der angekündigte Aha-Moment:
Die endliche Geometrie hat den Verzweigungsschnitt zum Verschwinden
gebracht, und $u_0$ ist eine meromorphe Funktion, deren einzige
Singularitäten isolierte Pole sind.

Diese Pole sind die Nullstellen des Nenners.
Der komplexe Sinus Hyperbolicus hängt über die Identität
$\sinh(z) = -i \sin(iz)$ mit dem Sinus zusammen; seine Nullstellen liegen
deshalb sämtlich auf der imaginären Achse, bei $z = i\nu\pi$ mit
ganzzahligem $\nu$.
Aus $\sinh(\pi\sqrt{s}) = 0$ folgt somit $\pi\sqrt{s} = i\nu\pi$, also
\begin{equation}
    s_\nu
    =
    -\nu^2,
    \qquad
    \nu = 1, 2, 3, \dots,
    \label{laplace:eq:polstellen}
\end{equation}
wobei $s = 0$ ausscheidet, weil dort auch der Zähler verschwindet und die
Singularität hebbar ist.
Die Bildfunktion besitzt somit \emph{unendlich viele} einfache Pole, die
sich mit quadratisch wachsenden Abständen auf der negativen reellen Achse
aufreihen, wie Abbildung~\ref{laplace:fig:pol_kette} zeigt.

\begin{figure}
    \centering
    % ============================================================================
% Bild 4: Bemasste "technische Zeichnung" der unendlichen Pol-Kette
% ============================================================================
% MATHEMATISCHER HINTERGRUND / DESIGN-ENTSCHEIDUNGEN:
%
% * Dargestellt werden die Pole s_nu = -nu^2 (nu = 1, 2, 3, ...) der
%   Bildfunktion u_0(x,s) = sinh((pi-x)sqrt(s))/sinh(pi sqrt(s)) des
%   Stabes der Laenge l = pi, sowie die Bogenfolge C_n aus dem Paper
%   mit den im URSPRUNG zentrierten Radien R_n = (n + 1/2)^2, die
%   GENAU ZWISCHEN den Polen hindurchlaufen:
%       R_1 = 2.25   (zwischen -1 und -4)
%       R_2 = 6.25   (zwischen -4 und -9)
%
% * MASSSTAB: 0.4 cm pro s-Einheit (linear bis zur Bruchstelle).
%   Damit liegen die Pole bei x = -0.4, -1.6, -3.6 cm und die
%   Bogenradien betragen 0.9 cm bzw. 2.5 cm.
%
% * BOGEN-GEOMETRIE: Die Boegen schliessen an die Bromwich-Gerade
%   Re(s) = gamma an (gamma = 1 s-Einheit = 0.4 cm), d.h. die Gerade
%   ist SEHNE der Kreise (konsistent zur D-Kontur-Abbildung).
%   Startwinkel aus cos(theta) = 0.4/R:
%       C_1: R = 0.9  -> theta = 63.61 deg, Endpunkt-y = 0.9*sin = 0.806
%       C_2: R = 2.5  -> theta = 80.79 deg, Endpunkt-y = 2.5*sin = 2.468
%
% * BEMASSUNG: Massketten unterhalb der Achse zeigen die Abstaende
%   1, 3, 5 (in s-Einheiten) zwischen 0,-1,-4,-9 -- die ungeraden
%   Zahlen machen das quadratische Wachstum |s_nu| = nu^2 sichtbar.
%   Technischer Stil: duenne graue Hilfslinien + Masslinien mit
%   Endstrichen (Bar-Tips statt Pfeilen, weil das kuerzeste Mass mit
%   0.4 cm zu schmal fuer beidseitige Pfeilspitzen ist).
%
% * BRUCHSTELLE: Doppelter Zickzack-Bruch (techn. Zeichnungsnorm)
%   nach -9; dahinter ein letzter Pol "-n^2" und ein Achsenstummel
%   mit Pfeil nach links (Fortsetzung der Kette bis -unendlich).
%
% * Bei s = 0 KEIN Pol (hebbare Singularitaet, Zaehler verschwindet
%   mit): offener Kreis statt Kreuz, graue Annotation "hebbar".
% ============================================================================
\scalebox{0.85}{
\begin{tikzpicture}

% ==============================
% 1. FARBEN DEFINIEREN
% ==============================
\definecolor{plotblue}{RGB}{0, 0, 200}
\definecolor{plotred}{RGB}{220, 30, 30}
\definecolor{plotgray}{RGB}{140, 140, 140}
\definecolor{plotbg}{RGB}{244, 244, 252}

% ==============================
% 2. BOEGEN C_n (zuunterst, damit alles andere darueber liegt)
% ==============================
% C_1: R = 0.9 cm, Sehne bei x = 0.4  ->  arc von 63.61 bis 296.39 Grad
\draw[thick, plotgray, dashed] (0.4, 0.806) arc (63.61:296.39:0.9);
% C_2: R = 2.5 cm, Sehne bei x = 0.4  ->  arc von 80.79 bis 279.21 Grad
\draw[thick, plotgray, dashed] (0.4, 2.468) arc (80.79:279.21:2.5);

% Bogen-Labels (jeweils knapp ausserhalb des Bogens, oben links)
\node[plotgray, font=\normalsize] at (-0.55, 1.15) {$C_1$};
\node[plotgray, font=\normalsize] at (-1.75, 2.15) {$C_2$};

% ==============================
% 3. BROMWICH-GERADE (Sehne der Boegen bei Re(s) = gamma)
% ==============================
\draw[thick, plotblue, decoration={markings, mark=at position 0.6 with {\arrow{Stealth[length=3mm, width=2mm]}}}, postaction={decorate}]
    (0.4, -2.468) -- (0.4, 2.468);
\draw[thick] (0.4, 0.1) -- (0.4, -0.1) node[below right, inner sep=2pt] {$\boldsymbol{\gamma}$};

% ==============================
% 4. ACHSEN MIT BRUCHSTELLE
% ==============================
% Hauptachse (rechts der Bruchstelle)
\draw[thick, -{Stealth[length=3.5mm, width=2.5mm]}] (-3.9, 0) -- (2.3, 0) node[right] {\textbf{Re}$(s)$};
% Achsenstummel links der Bruchstelle, Pfeil nach links (bis -unendlich)
\draw[thick, -{Stealth[length=3.5mm, width=2.5mm]}] (-4.34, 0) -- (-6.05, 0);
% Doppelter Zickzack-Bruch (techn. Norm fuer "hier fehlt ein Stueck")
\draw[thick] (-3.92, -0.3) -- (-4.04, -0.06) -- (-3.96, 0.06) -- (-4.08, 0.3);
\draw[thick] (-4.14, -0.3) -- (-4.26, -0.06) -- (-4.18, 0.06) -- (-4.30, 0.3);
% Kurze Im-Achse
\draw[thick, -{Stealth[length=3.5mm, width=2.5mm]}] (0, -2.9) -- (0, 3.05) node[above] {\textbf{Im}$(s)$};

% ==============================
% 5. POLE (Kreuze) UND HEBBARE STELLE
% ==============================
\newcommand{\pole}[3][plotred]{
    \draw[thick, #1, line cap=round] (#2-0.12,#3-0.12) -- (#2+0.12,#3+0.12);
    \draw[thick, #1, line cap=round] (#2-0.12,#3+0.12) -- (#2+0.12,#3-0.12);
}

% Pole bei s = -1, -4, -9 (Massstab 0.4 cm/Einheit) und -n^2
\pole{-0.4}{0}
\pole{-1.6}{0}
\pole{-3.6}{0}
\pole{-5.0}{0}

% Pol-Beschriftungen OBERHALB der Achse (unten laeuft die Masskette)
\node[anchor=south, font=\small] at (-0.4, 0.22) {$-1$};
\node[anchor=south, font=\small] at (-1.6, 0.22) {$-4$};
\node[anchor=south, font=\small] at (-3.6, 0.22) {$-9$};
\node[anchor=south, font=\small] at (-5.0, 0.22) {$-n^2$};

% Hebbare Singularitaet bei s = 0: roter Kreis mit rotem Kreuz
\draw[thick, plotred] (0,0) circle (3.0pt);
\draw[thick, plotred, line cap=round] (-0.075,-0.075) -- (0.075,0.075);
\draw[thick, plotred, line cap=round] (-0.075,0.075) -- (0.075,-0.075);



% ==============================
% 7. LEGENDE
% ==============================
\node[draw, thick, rounded corners=3pt, fill=white, inner xsep=5pt, inner ysep=4pt, anchor=north west] at (1.3, 3.05) {
    \renewcommand{\arraystretch}{1.2}
    \begin{tabular}{@{}c l@{}}

        \tikz[baseline=-0.6ex]{
            \draw[thick, plotred, line cap=round] (-0.1,-0.1) -- (0.1,0.1) (-0.1,0.1) -- (0.1,-0.1);
        } & Polstelle \\

        \tikz[baseline=-0.6ex]{
            \draw[thick, plotred] (0,0) circle (3.0pt);
            \draw[thick, plotred, line cap=round] (-0.075,-0.075) -- (0.075,0.075);
            \draw[thick, plotred, line cap=round] (-0.075,0.075) -- (0.075,-0.075);
        } & Hebbare Singularit\"at \\

        \tikz[baseline=-0.6ex]{
            \draw[thick, plotblue, ->-=0.6] (-0.3, 0) -- (0.3, 0);
        } & Bromwich-Pfad \\

        \tikz[baseline=-0.6ex]{
            \draw[thick, plotgray, dashed] (-0.3, 0) -- (0.3, 0);
        } & Bogen $C_n$ \\

    \end{tabular}
};

\end{tikzpicture}
}

    \caption{Die Singularitäten des Stabes endlicher Länge:
    unendlich viele einfache Pole bei $s_\nu = -\nu^2$ auf der negativen
    reellen Achse.
    Die Abstände wachsen quadratisch, die Wellenlinie deutet die
    Fortsetzung der Kette bis zum $n$-ten Pol und darüber hinaus an.
    Die gestrichelten Bögen $C_n$ verlaufen jeweils zwischen zwei Polen
    und schliessen den Bromwich-Pfad für die Residuenauswertung.}
    \label{laplace:fig:pol_kette}
\end{figure}

Die Residuen an diesen Polen berechnen wir mit der Quotientenregel:
Hat $f_2$ bei $s_0$ eine einfache Nullstelle und ist $f_1(s_0) \neq 0$, so
ist das Residuum von $\frac{f_1}{f_2}$ dort gleich
$\frac{f_1(s_0)}{f_2'(s_0)}$.
Mit $f_2(s) = \sinh(\pi\sqrt{s})$ und
$f_2'(s) = \frac{\pi}{2\sqrt{s}}\cosh(\pi\sqrt{s})$ ergibt sich nach
kurzer Rechnung mit $\sqrt{s_\nu} = i\nu$ das erfreulich einfache Resultat
\begin{equation}
    \operatorname{Res}\bigl(u_0(x,s), -\nu^2\bigr)
    =
    \frac{2}{\pi} \, \nu \sin(\nu x).
    \label{laplace:eq:stab_residuen}
\end{equation}

Für die Inversion können wir die D-Kontur nicht mehr in einem Zug
schliessen, denn jeder feste Kreisbogen würde nur endlich viele der
unendlich vielen Pole erfassen, und ein Bogen durch einen Pol hindurch ist
verboten.
Stattdessen schliessen wir den Bromwich-Pfad durch eine \emph{Folge} von
wiederum im Ursprung zentrierten Bögen $C_n$, deren Radien
\begin{equation}
    R_n
    =
    \Bigl(n + \tfrac{1}{2}\Bigr)^2
    \label{laplace:eq:bogenradien}
\end{equation}
zwischen den Polen $-n^2$ und $-(n+1)^2$ hindurchlaufen.
Exakt mittig liegt der Radius dabei nur im transformierten
$\sqrt{s}$-Raum, wo $\sqrt{R_n} = n + \tfrac{1}{2}$ die Nullstellen $n$
und $n+1$ halbiert; in der $s$-Ebene selbst liegt
$R_n = n^2 + n + \tfrac{1}{4}$ etwas näher am inneren Pol, was für das
Argument aber keine Rolle spielt.
Der $n$-te Bogen schliesst die ersten $n$ Pole ein:
Für jedes feste $n$ liefert der Residuensatz auf der geschlossenen Kurve
aus Geradenstück und Bogen $C_n$ die endliche Partialsumme der ersten $n$
Residuen.
Es wird also nichts von den Bögen \emph{subtrahiert} -- ihr Beitrag geht
für $n \to \infty$ schlicht gegen null, und die Partialsummen konvergieren
gegen die unendliche Residuenreihe.
Für die Anregung mit einem Wärmeimpuls $A_0(t) = \delta(t)$, also
$a_0(s) = 1$, erhalten wir so
\begin{equation}
    U_0(x,t)
    =
    \lim_{n \to \infty}
    \frac{2}{\pi}
    \sum_{\nu=1}^{n}
    \nu \sin(\nu x) \, e^{-\nu^2 t}
    =
    \frac{2}{\pi}
    \sum_{\nu=1}^{\infty}
    \nu \sin(\nu x) \, e^{-\nu^2 t}.
    \label{laplace:eq:stab_loesung}
\end{equation}
Ein Wort zur Sorgfalt ist hier angebracht.
Der Nachweis, dass die Bogenintegrale tatsächlich verschwinden, ist der
entscheidende und in technischen Texten am häufigsten unterschlagene
Schritt; ohne ihn hat die Reihe nur heuristischen Charakter, wie
\cite{laplace:doetsch1961} mit Nachdruck betont.
Das Lemma von Jordan aus
Abschnitt~\ref{laplace:subsection:residuensatz} greift hier nämlich nicht
unmittelbar:
Es verlangt, dass die Bildfunktion für $|s| \to \infty$
\emph{gleichmässig} gegen null strebt, doch wegen der unendlich vielen
Pole kommt der Nenner $\sinh(\pi\sqrt{s})$ in jeder noch so weit
draussen liegenden Umgebung der negativen reellen Achse seinen
Nullstellen beliebig nahe -- von gleichmässigem Abfall kann keine Rede
sein.
Der Nachweis muss deshalb gezielt auf den Radien $R_n$ geführt werden,
die den Polen ausweichen und auf denen $|\sinh(\pi\sqrt{s})|$ nach unten
beschränkt bleibt.
Für die hier auftretenden Bildfunktionen lässt sich der Nachweis führen,
und die entsprechenden Abschätzungen auf den grossen Kreisbögen werden in
\cite{laplace:carslaw1959} für eine ganze Klasse von
Wärmeleitungsproblemen sauber ausgearbeitet.

Das Ergebnis~\eqref{laplace:eq:stab_loesung} verdient einen zweiten Blick,
denn wir kennen es bereits aus einem ganz anderen Kontext.
Es ist exakt die Fourier-Reihe, die die klassische Separation der
Variablen liefert:
Eigenmoden $\sin(\nu x)$ des endlichen Stabes, die mit den Raten
$e^{-\nu^2 t}$ abklingen.
Die Pole der Bildfunktion sind nichts anderes als die Eigenwerte des
räumlichen Problems, und die Residuen sind die zugehörigen Eigenmoden.
Der Residuensatz hat die Modalzerlegung nicht vorausgesetzt, sondern
selbstständig wiederentdeckt.

Damit klärt sich auch der Unterschied zum halbunendlichen Stab endgültig.
Für allgemeine Länge $l$ liegen die Pole bei
$s_\nu = -\frac{\nu^2\pi^2}{l^2}$, ihr Abstand ist also von der Ordnung
$\frac{\pi^2}{l^2}$.
Lassen wir $l$ wachsen, so rücken die Pole immer dichter zusammen, und im
Grenzfall $l \to \infty$ verschmilzt die diskrete Polkette zu einem
Kontinuum auf der negativen reellen Achse, eben dem Verzweigungsschnitt
aus Abschnitt~\ref{laplace:subsection:verzweigung}.
Gleichzeitig geht die Residuenreihe über die Raten $e^{-\nu^2\pi^2 t/l^2}$
in das Schnittintegral über die Raten $e^{-rt}$ aus
Gleichung~\eqref{laplace:eq:keyhole_loesung} über.
Diskretes Spektrum bei endlicher Geometrie, kontinuierliches Spektrum bei
unendlicher Geometrie:
Die komplexe $s$-Ebene bildet die Physik des Problems mit bestechender
Klarheit ab.

%
% Teil 5: Numerische Inversion -- Oszillationsproblem des
% Bromwich-Integrals, Pfaddeformation nach Cauchy, Talbot-Kontur,
% Trapezregel und Konvergenzvergleich
%
\section{Numerische Inversion: Talbots Weg ins Tal des Integranden}
\label{laplace:subsection:talbot}

In der Praxis endet die analytische Inversion schnell in einer Sackgasse.
Viele Bildfunktionen besitzen keine geschlossen darstellbare
Originalfunktion, oder die gefundenen Reihen und Integrale konvergieren
für die interessierenden Zeiten quälend langsam.
Gefragt ist dann eine numerische Auswertung des Bromwich-Integrals, und
der naheliegendste Versuch wäre eine \emph{Quadratur}, also eine
numerische Integration über gewichtete Funktionswerte an endlich vielen
Stützstellen, direkt entlang der Bromwich-Geraden.
Dieser Versuch scheitert spektakulär.
Auf der Geraden $s = \gamma + i\omega$ gilt
$e^{st} = e^{\gamma t} \left( \cos \omega t + i \sin \omega t \right)$;
der Betrag des Exponentialfaktors ist konstant, und der Integrand
oszilliert unablässig, während er nur so langsam abklingt, wie es die
Bildfunktion $F(s)$ selbst vorgibt.
Die Quadratur muss dann riesige, sich fast vollständig auslöschende
Beiträge aufsummieren, und die Genauigkeit wächst mit der Zahl der
Stützstellen nur schleppend.
Abbildung~\ref{laplace:fig:integrand_vergleich} zeigt dieses Verhalten für
das einfachste denkbare Beispiel $F(s) = \frac{1}{s}$.

\begin{figure}
    \centering
    % ============================================================================
% PLATZHALTER -- Bild 5: Entrollter Integrand, Bromwich vs. Talbot
% ============================================================================
% REGIEANWEISUNG FUER DIE GRAFIK-TOOLCHAIN (bevorzugt matplotlib/pgfplots,
% als PDF eingebunden oder direkt pgfplots; Schrift Times 10pt wie Text):
%
% DATENGRUNDLAGE (echte Berechnung, F(s)=1/s, t=1):
% - BROMWICH: s = gamma + i*omega mit gamma=1; Integrand
%   Re[ e^{st} F(s) ] = Re[ e^{t(gamma+i*omega)} / (gamma+i*omega) ]
%   als Funktion des Kurvenparameters omega in [-40, 40].
%   -> oszillierende Kurve mit nur ~1/|omega| abfallender Amplitude.
% - TALBOT: Kontur s(theta) = sigma + lambda(theta*cot(theta) + i*beta*theta)
%   mit z.B. sigma=0, lambda=2, beta=1, theta in (-pi, pi);
%   (beta ist der Formparameter, im Paper-Text s_beta; in der
%   Originalliteratur oft nu genannt -- hier NICHT nu beschriften!)
%   Integrand Re[ e^{t s(theta)} F(s(theta)) s'(theta) / (2*pi*i) ]
%   als Funktion von theta.
%   -> glatte, glockenfoermige Kurve, an den Raendern theta -> +-pi
%      praktisch null (superexponentieller Abfall).
%
% DARSTELLUNG:
% - ZWEI Panels nebeneinander (gleiche Hoehe), oder ein Panel mit zwei
%   Kurven, falls die Skalen vertraeglich sind; bevorzugt ZWEI PANELS:
%   * links:  "Bromwich-Gerade", Kurve in ROT (plotred), x-Achse
%     "$\omega$", y-Achse "Integrand".
%   * rechts: "Talbot-Kontur", Kurve in BLAU (plotblue), x-Achse
%     "$\theta$", y-Achse identisch skaliert beschriftet.
% - Beide y-Achsen mit derselben Einheit; falls die Amplituden stark
%   differieren, y-Limits pro Panel waehlen, aber im Caption-Text bleibt
%   die Aussage qualitativ.
% - Dezentes Gitter (grau, duenn), keine Titelzeile im Bild (Caption
%   uebernimmt das), Legenden nur falls beide Kurven in einem Panel.
% - Auf der Talbot-Kurve zusaetzlich N=12 Stuetzstellen der Trapezregel
%   als Punkte markieren (blaue Punkte), um zu zeigen, wie wenige
%   Auswertungen genuegen.
%
% GROESSE: Gesamtbreite ca. 12cm, Hoehe ca. 4.5cm.
% ============================================================================
\scalebox{0.8}{
\begin{tikzpicture}

% ==============================
% 1. FARBEN DEFINIEREN
% ==============================
\definecolor{plotblue}{RGB}{0, 0, 200}
\definecolor{plotred}{RGB}{220, 30, 30}
\definecolor{plotgray}{RGB}{140, 140, 140}
\definecolor{plotbg}{RGB}{244, 244, 252}

% ==============================
% 2. PGFPLOTS SETUP
% ==============================
\pgfplotsset{table/search path={.,Latex/images,images,papers/laplace/images}}

% ==============================
% 3. BROMWICH PANEL (Links)
% ==============================
\begin{axis}[
    name=ax1,
    width=7cm,
    height=5.5cm,
    grid=major,
    grid style={plotgray!30, thin},
    axis lines=left,
    axis line style={thick},
    tick label style={font=\small},
    label style={font=\normalsize},
    enlarge x limits=false,
    xlabel={$\omega$},
    ylabel={Integrand},
    xmin=-40, xmax=40,
    ymin=-0.8, ymax=3.0,
    title={\textbf{Bromwich-Gerade}},
    title style={font=\normalsize, yshift=-1.5ex}
]
\addplot[
    thick, 
    plotred,
    ] table[x=omega, y=integrand, col sep=comma] {papers/laplace/images/data_integrand_bromwich.csv};
\end{axis}

% ==============================
% 3. TALBOT PANEL (Rechts)
% ==============================
\begin{axis}[
    name=ax2,
    at={(ax1.outer east)},
    anchor=outer west,
    xshift=1.5cm,
    width=7cm,
    height=5.5cm,
    grid=major,
    grid style={plotgray!30, thin},
    axis lines=left,
    axis line style={thick},
    tick label style={font=\small},
    label style={font=\normalsize},
    enlarge x limits=false,
    xlabel={$\theta$},
    xmin=-3.1415, xmax=3.1415,
    ymin=-0.3, ymax=1.3,
    xtick={-3.1415, -1.5708, 0, 1.5708, 3.1415},
    xticklabels={$-\pi$, $-\pi/2$, $0$, $\pi/2$, $\pi$},
    title={\textbf{Talbot-Kontur}},
    title style={font=\normalsize, yshift=-1.5ex}
]
\addplot[
    thick, 
    plotblue,
    ] table[x=theta, y=integrand, col sep=comma] {papers/laplace/images/data_integrand_talbot.csv};

\addplot[
    only marks,
    mark=*,
    mark size=1.5pt,
    plotblue,
    ] table[x=theta, y=integrand, col sep=comma] {papers/laplace/images/data_integrand_talbot_pts.csv};
\end{axis}

\end{tikzpicture}
}

    \caption{Der entlang der Bogenlänge entrollte Realteil des
    Inversionsintegranden für $F(s) = \frac{1}{s}$ bei festem $t$.
    Auf der Bromwich-Geraden (rot) oszilliert die Funktion mit nur
    langsam abfallender Amplitude, auf der Talbot-Kontur (blau) ist sie
    eine glatte, rasch abklingende Kurve, die sich mit wenigen
    Stützstellen präzise integrieren lässt.}
    \label{laplace:fig:integrand_vergleich}
\end{figure}

Die Rettung liegt in einer Freiheit, die wir längst besitzen, aber bisher
nur zum Schliessen der Kontur genutzt haben.
Nach dem Cauchyschen Integralsatz dürfen wir den Integrationsweg beliebig
deformieren, solange wir dabei keine Singularität von $F(s)$ überstreichen
und $F(s)$ im überstrichenen Gebiet gleichmässig gegen null abfällt; der
Wert des Integrals bleibt dabei exakt gleich.
Diese Idee ist alt, denn schon \cite{laplace:carslaw1959} verformt den
Bromwich-Pfad zu Kurven, die links im Unendlichen beginnen und enden, um
die Konvergenz der Integrale nachzuweisen.
Talbot hat daraus in seiner wegweisenden Arbeit
\cite{laplace:talbot1979} ein numerisches Verfahren gemacht.
Sein Gedanke:
Man biege die beiden Enden des Pfades in die linke Halbebene, sodass
$\text{Re}(s) \to -\infty$ läuft.
Dort fällt der Faktor $|e^{st}| = e^{t\,\text{Re}(s)}$ exponentiell ab, der
Integrand wird winzig, und die Oszillationen verlieren jede Bedeutung.
Anschaulich verlässt der Pfad den Bergkamm der Betragsfläche
$|e^{st}F(s)|$ und steigt ins \emph{Tal} der linken Halbebene hinab.
In der logarithmischen Farbskala von
Abbildung~\ref{laplace:fig:betrags_landschaft} wird dieser Abstieg
unmittelbar sichtbar:
Die Bromwich-Gerade verharrt in den hellen Zonen konstanter Höhe,
während die Talbot-Kontur beidseitig in die dunklen Bereiche des
exponentiell kleinen Betrags eintaucht.

\begin{figure}
    \centering
    % ============================================================================
% PLATZHALTER -- Bild 6: Heatmap |e^{st} F(s)| mit Bromwich- und
%                        Talbot-Pfad
% ============================================================================
% REGIEANWEISUNG FUER DIE GRAFIK-TOOLCHAIN (matplotlib empfohlen,
% Export als PDF; Schrift Times 10pt; KEINE transparenten PNG-Teile):
%
% DATENGRUNDLAGE (F(s)=1/s, t=1):
% - Gitter der s-Ebene: Re(s) in [-15, 5], Im(s) in [-12, 12],
%   Aufloesung >= 800x800.
% - Farbwert: log10( |e^{s*t} / s| ), geclippt auf z.B. [-8, 2].
%
% DARSTELLUNG:
% - 2D-CONTOUR-HEATMAP (pcolormesh oder contourf mit vielen Stufen),
%   sequentielle Colormap (z.B. viridis); COLORBAR rechts mit
%   Beschriftung "$\log_{10}|e^{st}F(s)|$".
% - Zusaetzlich duenne weisse Konturlinien bei ganzzahligen log10-Stufen,
%   damit die "Landschaft" (Bergkamm rechts, Tal links) plastisch wird.
% - SINGULARITAET s=0: kleines weisses/rotes Kreuz mit Label "Pol".
% - BROMWICH-PFAD: vertikale ROTE Linie bei Re(s)=1 ueber die volle
%   Bildhoehe, Pfeilspitze nach oben, Label "$\Gamma_B$" (rot).
% - TALBOT-PFAD: BLAUE Kurve s(theta)=sigma+lambda(theta*cot(theta)
%   + i*beta*theta) mit denselben Parametern wie in Bild 5 (sigma=0,
%   lambda=2, beta=1; Formparameter beta, nicht nu beschriften),
%   Pfeilspitzen in Durchlaufrichtung (von unten nach
%   oben), Label "$\Gamma_T$" (blau).
%   Die Kurve muss erkennbar ins dunkle "Tal" links abtauchen und den
%   Pol rechts umschliessen.
% - Textannotation "Tal: $|e^{st}|$ exponentiell klein" im linken
%   Bildbereich (weiss, small); optional "Bergkamm" rechts.
% - Achsenbeschriftung "Re(s)", "Im(s)"; kein Bildtitel.
%
% ALTERNATIVE (falls 3D gewuenscht): 3D-Surface derselben Daten mit
% Blick von schraeg oben, Pfade auf die Flaeche projiziert; 2D-Version
% ist fuer den Druck aber besser lesbar und wird bevorzugt.
%
% GROESSE: ca. 11cm x 8cm inkl. Colorbar, 300dpi falls Rasterexport.
% ============================================================================
\scalebox{0.85}{
\begin{tikzpicture}

% ==============================
% 1. FARBEN DEFINIEREN
% ==============================
\definecolor{plotblue}{RGB}{0, 0, 200}
\definecolor{plotred}{RGB}{220, 30, 30}
\definecolor{plotgray}{RGB}{140, 140, 140}
\definecolor{plotbg}{RGB}{244, 244, 252}

% PGFPLOTS SETUP
\pgfplotsset{table/search path={.,Latex/images,images,papers/laplace/images}}

\begin{axis}[
    width=11cm,
    height=9.5cm,
    enlargelimits=false,
    axis on top,
    xlabel={$\operatorname{Re}(s)$},
    ylabel={$\operatorname{Im}(s)$},
    xmin=-15, xmax=5,
    ymin=-12, ymax=12,
    tick label style={font=\small},
    label style={font=\normalsize},
    colorbar,
    colormap/viridis,
    point meta min=-8,
    point meta max=2,
    colorbar style={
        ylabel={$\log_{10}|e^{st}F(s)|$},
        ytick={-8,-6,-4,-2,0,2},
        ylabel style={font=\normalsize},
        tick label style={font=\small}
    }
]

% Hintergrund-Heatmap (Bild, darauf projiziert)
\addplot graphics [xmin=-15, xmax=5, ymin=-12, ymax=12] {papers/laplace/images/heatmap_bg.png};

% Singularitaet s=0 (Pol)
\draw[thick, white, line cap=round] (axis cs:-0.4,-0.4) -- (axis cs:0.4,0.4) (axis cs:-0.4,0.4) -- (axis cs:0.4,-0.4);
\draw[thick, plotred, line cap=round] (axis cs:-0.2,-0.2) -- (axis cs:0.2,0.2) (axis cs:-0.2,0.2) -- (axis cs:0.2,-0.2);
\node[white, left, font=\small, xshift=-2pt] at (axis cs:-0.2, 0) {\textbf{Pol}};

% Bromwich-Gerade (rot)
\draw[thick, plotred, decoration={markings, mark=at position 0.65 with {\arrow{Stealth[length=3mm, width=2mm]}}}, postaction={decorate}] 
    (axis cs:1, -12) -- (axis cs:1, 12);
\node[plotred, right, font=\large] at (axis cs:1.2, 9) {$\Gamma_B$};

% Talbot-Kontur (blau)
\addplot[
    thick, plotblue,
    decoration={markings, mark=at position 0.35 with {\arrow{Stealth[length=3mm, width=2mm]}}},
    decoration={markings, mark=at position 0.65 with {\arrow{Stealth[length=3mm, width=2mm]}}},
    postaction={decorate}
    ] table[col sep=comma, x=re, y=im] {papers/laplace/images/talbot_contour_heatmap.csv};
\node[plotblue, left, font=\large] at (axis cs:-3, 6) {$\Gamma_T$};

% Annotation Tal removed

% ==============================
% 2. LEGENDE
% ==============================
\filldraw[fill=white, draw=black, thick, rounded corners=3pt] (axis cs:-14.5, 11.5) rectangle (axis cs:-2.5, 8.5);
\draw[thick, plotred, decoration={markings, mark=at position 0.55 with {\arrow{Stealth[length=3mm, width=2mm]}}}, postaction={decorate}] (axis cs:-14.0, 10.5) -- (axis cs:-12.5, 10.5);
\node[right] at (axis cs:-12.0, 10.5) {Bromwich-Gerade $\Gamma_B$};
\draw[thick, plotblue, decoration={markings, mark=at position 0.55 with {\arrow{Stealth[length=3mm, width=2mm]}}}, postaction={decorate}] (axis cs:-14.0, 9.5) -- (axis cs:-12.5, 9.5);
\node[right] at (axis cs:-12.0, 9.5) {Talbot-Kontur $\Gamma_T$};

\end{axis}
\end{tikzpicture}
}

    \caption{Betragslandschaft $|e^{st}F(s)|$ über der komplexen
    $s$-Ebene für $F(s) = \frac{1}{s}$ zur Zeit $t = 1$ in
    logarithmischer Farbskala.
    Die Bromwich-Gerade (rot) verläuft auf konstanter Höhe durch das
    oszillierende Gebiet, während die Talbot-Kontur (blau) die
    Singularität umschliesst und beidseitig ins exponentiell tiefe Tal
    der linken Halbebene abtaucht.}
    \label{laplace:fig:betrags_landschaft}
\end{figure}

Konkret verwendet Talbot die Kontur
\begin{equation}
    s(\theta)
    =
    \sigma + \lambda \, s_{\beta}(\theta),
    \qquad
    s_{\beta}(\theta)
    =
    \theta \cot \theta + i \beta \theta,
    \qquad
    -\pi < \theta < \pi,
    \label{laplace:eq:talbot_kontur}
\end{equation}
mit reellen Parametern $\lambda > 0$, $\beta > 0$ und einer Verschiebung
$\sigma$; für den Formparameter schreibt die Originalliteratur meist
$\nu$, was wir hier vermeiden, weil dieser Buchstabe in
Abschnitt~\ref{laplace:subsection:endlicher_stab} bereits als Polindex
vergeben ist.
Die Verschiebung $\sigma$ übernimmt dabei die Rolle der
Bromwich-Konstanten $\gamma$ aus
Gleichung~\eqref{laplace:eq:bromwich_integral}; die Umbenennung folgt der
numerischen Literatur und bezeichnet denselben Rechtsversatz des
Integrationsweges.
Für $\theta \to \pm\pi$ strebt $\theta \cot\theta \to -\infty$, die Kurve
öffnet sich also wie gewünscht nach links und legt sich um die negative
reelle Achse herum; sämtliche Singularitäten unserer
Wärmeleitungsbeispiele, ob Polkette oder Verzweigungsschnitt, werden von
ihr umschlossen.
Setzt man die Parametrisierung in das Bromwich-Integral ein, so entsteht
das Integral über ein endliches Intervall
\begin{equation}
    f(t)
    =
    \frac{\lambda e^{\sigma t}}{2\pi i}
    \int_{-\pi}^{\pi}
    e^{\lambda t s_{\beta}(\theta)}
    \,
    F\bigl(\sigma + \lambda s_{\beta}(\theta)\bigr)
    \,
    s_{\beta}'(\theta)
    \, d\theta.
    \label{laplace:eq:talbot_integral}
\end{equation}
Der Integrand ist analytisch in $\theta$ und verschwindet an beiden
Intervallenden samt allen Ableitungen exponentiell.
Für genau solche Integranden ist die einfachste aller Quadraturformeln,
die Trapezregel, nicht etwa grob, sondern \emph{geometrisch konvergent}:
Jede zusätzliche Stützstelle multipliziert den Fehler mit einem festen
Faktor kleiner als eins \cite{laplace:trefethen2006}.
Der Grund dafür steckt in der Euler-Maclaurin-Formel, nach der der
Trapezfehler aus Randtermen mit den Ableitungen des Integranden an den
Intervallenden besteht:
Da unser Integrand bei $\theta = \pm\pi$ samt allen Ableitungen
abstirbt, verschwinden diese Randterme allesamt, und übrig bleibt nur
ein exponentiell kleiner Rest.
Mit $n$ äquidistanten Knoten $\theta_j = \frac{j\pi}{n}$ und unter
Ausnutzung der Symmetrie der Kontur zur reellen Achse ergibt sich die
Näherung
\begin{equation}
    \tilde{f}(t)
    =
    \frac{\lambda e^{\sigma t}}{n}
    \sum_{j=0}^{n-1}
    \operatorname{Re}
    \left[
        e^{\lambda t s_{\beta}(\theta_j)}
        \,
        F\bigl(\sigma + \lambda s_{\beta}(\theta_j)\bigr)
        \,
        \frac{s_{\beta}'(\theta_j)}{i}
    \right],
    \label{laplace:eq:talbot_summe}
\end{equation}
wobei der Summand für $j = 0$ mit dem Faktor $\frac{1}{2}$ gewichtet wird
\cite{laplace:murli1990}.
Der Randknoten $j = n$, also $\theta = \pi$, fehlt in der Summe nicht aus
Nachlässigkeit:
Dort strebt $\text{Re}\,s_{\beta}(\theta) \to -\infty$, und der Integrand
verschwindet exakt, sodass sein Beitrag zur Trapezsumme entfällt.
Der Realteil-Operator und die halbierten Summationsgrenzen sind dabei
kein Formalismus, sondern gezielte Ausnutzung der Kontursymmetrie:
Weil die Kontur spiegelsymmetrisch zur reellen Achse verläuft und für
reelles $f(t)$ die Beziehung $F(\bar{s}) = \overline{F(s)}$ gilt, liefern
die Stützstellen mit negativem $\theta_j$ exakt die konjugierten Beiträge
ihrer Spiegelpartner.
Statt über alle $2n$ Knoten des vollen Intervalls $(-\pi, \pi)$ zu
summieren, genügen daher die $n$ Knoten mit $\theta_j \geq 0$; die
Paarbildung steckt im Realteil, und der Rechenaufwand halbiert sich.
Die gesamte Inversion schrumpft damit auf eine Handvoll Auswertungen der
Bildfunktion an komplexen Stützstellen zusammen.

Die Parameter sind dabei nicht frei wählbar, sondern müssen zum Problem
passen.
Die Verschiebung $\sigma$ wird rechts der Realteile aller Singularitäten
gewählt, damit die Kontur die Singularitäten tatsächlich umschliesst.
Entscheidend ist die Skalierung von $\lambda$:
Damit die geometrische Konvergenz greift, muss die Kontur mit der
Stützstellenzahl wachsen und mit der Inversionszeit schrumpfen, also
$\lambda \propto \frac{N}{t}$, wobei $N = 2n$ die Gesamtzahl der
Stützstellen auf der vollen Kontur bezeichnet, von denen
Gleichung~\eqref{laplace:eq:talbot_summe} dank der Symmetrie nur die
Hälfte auswertet.
In der von Weideman optimierten Form lautet die Kontur beispielsweise
\begin{equation}
    s(\theta)
    =
    \frac{N}{t}
    \bigl(
        0{,}5017 \, \theta \cot(0{,}6407\,\theta)
        - 0{,}6122
        + 0{,}2645 \, i\theta
    \bigr),
    \label{laplace:eq:weideman_kontur}
\end{equation}
mit der die Rate $O(3{,}89^{-N})$ erreicht wird
\cite{laplace:weideman2006}.
Man beachte, dass hier die \emph{gesamte} Kontur -- Verschiebung und
Skalierung gleichermassen -- mit dem Faktor $\frac{N}{t}$ wächst.
Genau diese Kopplung von Konturgrösse, Stützstellenzahl und
Inversionszeit ist der Kern der Parameterwahl und der Garant für die
geometrische Konvergenz; mit einer festen, von $N$ und $t$ unabhängigen
Kontur ginge sie verloren.
Die Kopplung hat aber auch eine praktische Konsequenz:
Die Stützstellen hängen von $t$ ab, sodass für jeden neuen Zeitpunkt die
Bildfunktion an neuen Punkten ausgewertet werden muss
\cite{laplace:trefethen2006}.

Wie gut das funktioniert, ist verblüffend.
Bereits mit $n = 11$ Stützstellen erreicht Talbots Verfahren typischerweise
Fehler der Ordnung $10^{-6}$, mit $n = 20$ etwa $10^{-11}$ und mit
$n = 35$ in doppelter Genauigkeit rund $10^{-20}$
\cite{laplace:talbot1979}.
Die sorgfältige Software-Implementierung von Murli und Rizzardi bestätigt
diese Zahlen und zeigt in direkten Vergleichen, dass Fourier-basierte
Verfahren für dieselbe Genauigkeit ein Vielfaches an Funktionsauswertungen
benötigen \cite{laplace:murli1990}.
Die moderne Analyse von Trefethen, Weideman und Schmelzer ordnet die
Talbot-Kontur in die Familie der \emph{Hankel-Konturen} ein.
Der Begriff ist dabei ein topologischer Überbegriff:
Er bezeichnet jede Kurve, die aus dem dritten Quadranten von
$\text{Re}(s) = -\infty$ kommend die negative reelle Achse umschlingt
und im zweiten Quadranten wieder nach $\text{Re}(s) = -\infty$
verschwindet -- die Schlüsselloch-Kontur aus
Abschnitt~\ref{laplace:subsection:verzweigung} gehört ebenso dazu wie
Parabeln, Hyperbeln oder Talbots Kotangens-Kontur.
Innerhalb dieser Familie beziffert die Analyse die Konvergenzraten mit
optimierten Parametern:
Parabeln erreichen $O(2{,}85^{-N})$, Hyperbeln $O(3{,}20^{-N})$ und
Talbots Kotangens-Konturen $O(3{,}89^{-N})$ bei $N$ Stützstellen
\cite{laplace:trefethen2006}; die zugrunde liegende Parameteroptimierung
stammt aus \cite{laplace:weideman2006}.
Diese Raten sind keine Eigenheit des Testbeispiels $F(s) = \frac{1}{s}$,
sondern theoretische Schranken, die für die gesamte Klasse von
Bildfunktionen mit Singularitäten auf oder nahe der negativen reellen
Achse gelten -- also für alle unsere Wärmeleitungsprobleme, ob Polkette
oder Verzweigungsschnitt.
Mit $N = 32$ ist damit Maschinengenauigkeit erreicht.
Abbildung~\ref{laplace:fig:konvergenz_vergleich} stellt dieses
geometrische Konvergenzverhalten der langsamen Konvergenz auf der
Bromwich-Geraden gegenüber.

\begin{figure}
    \centering
    % ============================================================================
% Bild 7: Konvergenzvergleich (Fehler vs. Stuetzstellen), nativ in pgfplots
% ============================================================================
% DATEN: data_konvergenz.csv, erzeugt von generate_konvergenz_data.py
% (Testproblem F(s) = 1/s, exakt f(t) = 1, t = 1.5).
%
%   * BROMWICH: Mittelpunktsregel auf s = 1 + i*omega, |omega| <= 2*sqrt(N).
%     Der Abschneidefehler ~ e^{gamma t}/(pi*t*Omega) ~ N^(-1/2) dominiert
%     (der Aliasing-Anteil faellt wie e^{-(pi/2)*sqrt(N)} weg):
%     langsame, leicht oszillierende algebraische Konvergenz.
%   * TALBOT: Mittelpunktsregel auf der Weideman-Kontur (Trefethen/
%     Weideman/Schmelzer 2006, Gl. (3.3), skaliert mit N/t).
%     Geometrische Konvergenz O(3.89^-N) bis ~3e-14 bei N ~ 28; DANACH
%     steigt der Fehler wieder an: der Konturscheitel liegt bei
%     Re(st) ~ 0.171*N, also verstaerkt der Faktor e^{st} den Rundungs-
%     fehler wie eps*e^{0.171*N} (Talbots Fehlerterm E_r) -- bei N = 100
%     sind das ~5e-9. Dieses V-foermige Verhalten ist ECHT und wird
%     bewusst gezeigt (Annotation "Rundungsfehler").
%
% DESIGN-ENTSCHEIDUNGEN:
%   * table/search path deckt beide Build-Kontexte ab: Buch-Build
%     (cwd = Latex/, Pfad images/...) und test_render.py (cwd =
%     Projektwurzel, Pfad Latex/images/...).
%   * Referenzgerade 0.5*3.89^(-N) nur bis N = 26 gezeichnet -- ab N = 28
%     ist sie in den Daten auf 1e-16 gesaettigt und wuerde sonst als
%     horizontale Linie auf dem Boden weiterlaufen.
%   * Umlaute als \"u-Escapes (test_render.py liest die Datei ohne
%     explizites Encoding; rohe UTF-8-Bytes wuerden verstuemmelt).
%   * 17 Dekaden (1e-17..1e0), Ticks alle 4 Dekaden; Maschinengenauigkeit
%     als gepunktete Linie bei eps ~ 1e-16 knapp ueber dem Achsenboden.
% ============================================================================
\begin{tikzpicture}

% ==============================
% 1. FARBEN DEFINIEREN
% ==============================
\definecolor{plotblue}{RGB}{0, 0, 200}
\definecolor{plotred}{RGB}{220, 30, 30}
\definecolor{plotgray}{RGB}{140, 140, 140}
\definecolor{plotbg}{RGB}{244, 244, 252}

% CSV in beiden Build-Kontexten auffindbar machen.
% REIHENFOLGE WICHTIG: Latex/images VOR images, denn im Werkstatt-Root
% existiert ein aelterer images/-Ordner mit veralteten Daten, der sonst
% die frische CSV verschatten wuerde. Beim Buch-Build (cwd = Latex/)
% existiert Latex/images nicht und es greift automatisch images/.
\pgfplotsset{table/search path={.,Latex/images,images,papers/laplace/images}}

\begin{semilogyaxis}[
    width=10.5cm,
    height=7.5cm,
    xlabel={Anzahl der St\"utzstellen $N$},
    ylabel={Absoluter Fehler $|f_{\text{num}} - 1|$},
    xmin=0, xmax=104,
    ymin=1e-17, ymax=1e0,
    xtick={0,20,40,60,80,100},
    ytick={1e0,1e-4,1e-8,1e-12,1e-16},
    grid=both,
    grid style={line width=.1pt, draw=gray!20},
    major grid style={line width=.2pt, draw=gray!50},
    legend cell align={left},
    legend style={at={(0.97,0.82)}, anchor=north east,
        nodes={scale=0.9, transform shape}, draw=black!30,
        rounded corners=2pt, inner xsep=4pt, inner ysep=3pt},
    tick label style={font=\footnotesize},
    label style={font=\small}
]

    % Maschinengenauigkeit (eps ~ 1e-16)
    \draw[black, thick, densely dotted]
        (axis cs:0, 1e-16) -- (axis cs:104, 1e-16);
    \node[black, above right, font=\footnotesize]
        at (axis cs:2, 1.3e-16) {Maschinengenauigkeit};

    % Referenzgerade O(3.89^-N), nur bis N = 24: danach ist sie in den
    % Daten gesaettigt UND wuerde unten in den Schriftzug
    % "Maschinengenauigkeit" hineinlaufen
    \addplot[plotgray, thick, dashed, forget plot]
        table [col sep=comma, x=N, y=err_ref,
               restrict x to domain=4:24] {papers/laplace/images/data_konvergenz.csv};
    % Label rechts oberhalb der Referenzgeraden im leeren Keil zwischen
    % Bromwich-Kurve (oben) und Talbot-Abfall (unten); die fruehere
    % Position am V-Boden (27, 1.5e-15) kollidierte optisch mit den
    % blauen Quadraten bei N = 24-28 und dem Schriftzug
    % "Maschinengenauigkeit"
    \node[plotgray, anchor=west, font=\footnotesize]
        at (axis cs:19, 2e-9) {$\mathcal{O}(3{,}89^{-N})$};

    % Bromwich-Gerade: langsame algebraische Konvergenz (rot, Kreise)
    \addplot[
        thick, plotred,
        mark=o, mark options={solid, scale=0.8}
    ] table [col sep=comma, x=N, y=err_bromwich] {papers/laplace/images/data_konvergenz.csv};
    \addlegendentry{Bromwich-Gerade}

    % Talbot-Kontur: geometrischer Abfall, dann Rundungsfehler-Anstieg
    % (blau, Quadrate)
    \addplot[
        thick, plotblue,
        mark=square, mark options={solid, scale=0.8}
    ] table [col sep=comma, x=N, y=err_talbot] {papers/laplace/images/data_konvergenz.csv};
    \addlegendentry{Talbot-Kontur}

    % Annotation am wieder ansteigenden Ast: eps * e^{0.171 N}
    % (etwas oberhalb der Kurve, damit die blauen Marker bei N ~ 76-84
    % nicht in den Schriftzug laufen)
    \node[plotgray, anchor=south, font=\footnotesize]
        at (axis cs:70, 8e-10) {Rundungsfehler};

\end{semilogyaxis}
\end{tikzpicture}

    \caption{Konvergenzvergleich der numerischen Inversion für
    $F(s) = \frac{1}{s}$ bei festem $t$:
    absoluter Fehler gegen die Anzahl der Stützstellen $N$ in
    halblogarithmischer Darstellung.
    Die Quadratur auf der Bromwich-Geraden (rot) verbessert sich nur
    langsam, während die Trapezregel auf der Talbot-Kontur (blau)
    geometrisch mit der Rate $O(3{,}89^{-N})$ konvergiert und um
    $N \approx 30$ nahe an die Maschinengenauigkeit herankommt.
    Für noch grössere $N$ steigt der Fehler wieder leicht an, weil der
    Faktor $e^{st}$ am Konturscheitel den Rundungsfehler wie
    $\varepsilon \, e^{0{,}17 N}$ verstärkt.}
    \label{laplace:fig:konvergenz_vergleich}
\end{figure}

Zwei Einschränkungen gehören der Vollständigkeit halber dazu.
Erstens muss $|F(s)|$ für $|s| \to \infty$ gleichmässig abfallen, damit
die Pfaddeformation erlaubt ist.
Zweitens müssen die Imaginärteile aller Singularitäten beschränkt sein,
denn die Kontur kann nur Singularitäten umschliessen, die nicht beliebig
weit nach oben oder unten ausweichen \cite{laplace:talbot1979}.
Und selbst innerhalb dieser Grenzen kostet jede Abweichung von der
reellen Achse:
Liegen Singularitäten mit grossem Imaginärteil $q_d$ vor, muss sich die
Kontur weit öffnen, um sie zu umschliessen, und die benötigte
Stützstellenzahl wächst proportional zum Produkt $v = q_d \, t$
\cite{laplace:talbot1979}.
Da die Lage der Singularitäten $q_d$ vom Problem fest vorgegeben ist,
wächst $v$ mit fortschreitender physikalischer Zeit $t$ zwangsläufig an
-- bei Schwingungsproblemen verliert das Verfahren seine Effizienz also
nicht durch einen behebbaren Defekt, sondern durch die Zeit selbst.
Unsere Wärmeleitungsprobleme mit ihren Polen und Schnitten auf der
negativen reellen Achse sind damit der Idealfall des Verfahrens, während
etwa die Schwingungsgleichung mit ihrer unendlichen Polkette auf der
imaginären Achse ausserhalb seines Anwendungsbereichs liegt.
Innerhalb dieses Bereichs ist die Methode jedoch so leistungsfähig, dass
sie heute weit über die Laplace-Inversion hinaus eingesetzt wird, etwa zur
Berechnung von Matrix-Exponentialfunktionen und zur Zeitintegration
parabolischer partieller Differentialgleichungen
\cite{laplace:trefethen2006}.

Damit schliesst sich der Bogen dieser Arbeit.
Aus der Korrespondenztabelle wurde das Bromwich-Integral, aus dem Integral
eine Residuensumme, aus dem Umgang mit Verzweigungsschnitten und
Polketten ein Verständnis dafür, wie die Geometrie des physikalischen
Problems die Singularitäten in der $s$-Ebene formt.
Die Talbot-Kontur schliesslich zeigt, dass dieselbe Funktionentheorie, die
die analytische Inversion trägt, auch das effizienteste numerische
Verfahren begründet:
Man muss das Bromwich-Integral nicht dort auswerten, wo es definiert
wurde, sondern darf es dorthin verschieben, wo es am schnellsten
verschwindet.


\printbibliography[heading=subbibliography]
\end{refsection}
